% LTeX: enabled=false

\documentclass{article}
\usepackage[utf8]{inputenc}

\newcommand{\nirpdftitle}{Trace Formula Review}
\usepackage{import}
\inputfrom{..}{nir}

\pagestyle{contentpage}

\title{The Trace Formula for the Impatient}
\author{Nir Elber}
\date{Summer 2026}
\rhead{\textit{TRACE FORMULA SPEEDRUN}}
\lhead{}

\setcounter{tocdepth}{2}

\begin{document}

\maketitle

% \begin{abstract}
% 	This document collects a variety of definitions and results from Lie theory.\todo{quaternionic}
% \end{abstract}

\tableofcontents

\newpage

\section{Tools}

\subsection{The Adelic Quotient}
\begin{notation}
	Throughout, $F$ is a global field, and $V(F)$ is the set of places of $F$, and $v\in V(F)$ is a place.
\end{notation}
\begin{notation}
	Throughout, $G$ is an affine algebraic group, usually reductive.
\end{notation}
\begin{notation}
	If $F$ is a number field, define $A_G$ to be the identity component of $A(\RR)$, where $A\subseteq\op{Res}_{F/\QQ}Z_G$ is the largest split torus. If $F$ is a function field of characteristic $p$, let $Z$ be the largest split torus in $\op{Res}_{F/\FF_p(t)}Z_G$, and define $A_G\subseteq Z(\FF_p(t)_\infty)$ to be the inverse image of the subgroup $\left(t^\ZZ\right)^d$ under a choice of isomorphism $Z\cong\mathbb G_m^d$.
\end{notation}
\begin{notation}
	Set
	\[G(\AA_F)^1\coloneqq\bigcap_{\chi\in X^*(G)}\ker\left(\left|\cdot\right|\circ\chi\right).\]
\end{notation}
\begin{notation}
	Set $[G]\coloneqq A_GG(F)\backslash G(\AA_F)$.
\end{notation}
\begin{proposition}
	If $F$ is a number field, the natural map $G(F)\backslash G(\AA_F)^1\to[G]$ is a homeomorphism. If $F$ is a function field, then $A_GG(\AA_F)^1$ has finite index in $G(\AA_F)$.
\end{proposition}
\begin{theorem}
	The quotient $G(F)\backslash G(\AA_F)^1$ has finite volume. It is compact if and only if every $F$-split torus $T\subseteq G$ has $T_{\overline F}\subseteq R(G_{\overline F})$.
\end{theorem}
\begin{theorem}[Dixmier--Malliavin]
	Suppose $G$ is reductive. Then any $f\in C_c^\infty(A_G\backslash G(\AA_F))$ is a finite sum of convolutions of pairs of functions in $C_c^\infty(A_G\backslash G(\AA_F))$.
\end{theorem}

\subsection{Decompositions}
\begin{theorem}[Levi]
	Suppose $G$ is reductive, and let $P\subseteq G$ be a parabolic subgroup. Let $N\subseteq P$ be the unipotent radical. Then the short exact sequence
	\[1\to N\to P\to P/N\to1\]
	splits.
\end{theorem}
\begin{theorem}[Iwasawa]
	Suppose $G$ is reductive, and let $P\subseteq G$ be a parabolic subgroup. Then there is a maximal compact subgroup $K\subseteq G(\AA_F)$ for which
	\[G(\AA_F)=P(\AA_F)K.\]
\end{theorem}
\begin{notation}
	Suppose $G$ is reductive, and let $P\subseteq G$ be a parabolic subgroup with Levi decomposition $P=MN$. Define $\mf a_M^*\coloneqq X^*(M)_\RR$ and $\mf a_M\coloneqq\op{Hom}(X^*(M),\RR)$. Then define $H_M\colon M(\AA_F)\to\mf a_M$ via
	\[\langle H_M(m),\lambda\rangle\coloneqq\log\left|\lambda(m)\right|.\]
	Because $K\cap M(\AA_F)\subseteq\ker H_M$, the function $H_P\colon G(\AA_F)\to\mf a_M$ given by $H_P(nmk)\coloneqq H_M(m)$ is well-defined.
\end{notation}
\begin{remark}
	Suppose $F$ is a number field. Then the restriction $H_P\colon A_M\to\mf a_M$ is an isomorphism: it is an iterated product of the isomorphism $\log\colon\RR^+\to\RR$.
\end{remark}
\begin{definition}[Siegel set]
	Suppose $G$ is reducitve, and let $P\subseteq G$ be a minimal parabolic subgroup with Levi decomposition $P=MN$ and Iwasawa decomposition $G(\AA_F)=P(\AA_F)K$. Let $Z_0\subseteq Z_G$ and $T_0\subseteq G^{\mathrm{der}}$ be maximal split tori. Then a \textit{Siegel set} in $G(\AA_F)^1$ is a set of the form
	\[\mathfrak S(t)\coloneqq\Omega_N\Omega_MA_{T_0}(t)K,\]
	where $\Omega_N\subseteq N(\AA_F)$ and $\Omega_M\subseteq M(\AA_F)^1$ are compact, and
	\[A_{T_0}(t)\coloneqq\{x\in A_{T_0}:\left|\alpha(x)\right|_\infty\ge t\text{ for all }\alpha\in\Delta\}.\]
\end{definition}
\begin{theorem}[Reduction theory]
	Suppose $G$ is reductive. Then there is a Siegel set $\mf S(t)$ such that
	\[G(F)\mf S(t)=G(\AA_F)^1.\]
\end{theorem}
\begin{remark}
	For $x\in\mf S(t)$, the $s_x\in A_{T_0}(t)$ piece is actually independent of the choice decomposition. Indeed, simply invert $H_P\colon\mf S(t)\to\mf a_M$.
\end{remark}
\begin{theorem}[Cartan]
	Let $P_0\subseteq G$ be a minimal parabolic subgroup with maximal split torus $A_0\subseteq P_0$, and let $K_0\subseteq G(F_v)$ be a maximal compact subgroup from the Iwasawa decomposition. If $v$ is finite, then
	\[G(F_v)=\bigcup_{\lambda\in X_*(A_0)^+}K_0\varpi^\lambda K_0.\]
\end{theorem}
\begin{proof}[Sketch]
	In the split case, define $K_0\coloneqq G(\OO_v)$.
\end{proof}
\begin{theorem}[Iwahori]
	Suppose $v$ is finite and $G$ is reductive. Let $P_0\subseteq G$ be a minimal parabolic subgroup, and let $K_0\subseteq G(F_v)$ be a maximal compact open subgroup with Iwasawa decomposition $G(F_v)=P_0(F_v)K_0$. Let $P_0=MN^+$ be the Levi decomposition, and let $N^-$ consist of the negative root spaces. Then there is compact open neighborhood basis $\{K_n\}_{n\ge0}$ of $1\in G(F_v)$ such that the multiplication maps
	\[\left(N^+\cap K_n\right)\times(M\cap K_n)\times(N^-\cap K_n)\to K_n\]
	is an isomorphism for all $n$.
\end{theorem}
\begin{proof}[Sketch]
	In the split case, define $K_n$ to be the kernel of $G(\OO_v)\to G(\OO_v/\mf p^n)$. In general, pass to the split case by a Galois descent argument.
\end{proof}

\subsection{Functional Analysis}
\begin{definition}
	Let $V$ be a Hilbert space. For a bounded operator $A\colon V\to V$, define the \textit{adjoint} $A^*\colon V\to V$ to be the unique linear function such that
	\[\langle Av,w\rangle_V=\langle v,A^*w\rangle_V\]
	for all $v,w\in V$. It turns out that $A^*$ is a bounded operator.
\end{definition}
\begin{definition}
	Let $V$ be a Hilbert space. An operator $A\colon V\to V$ is \textit{positive} if and only if $\langle Av,v\rangle_V\ge0$ for all $v\in V$.
\end{definition}
\begin{lemma}
	Let $A$ be a positive bounded operator on a Hilbert space $V$. Then $A$ admits a unique positive square $\sqrt A$ for which $\sqrt A\circ\sqrt A=A$.
\end{lemma}
\begin{notation}
	Let $A$ be a bounded operator on a Hilbert space $V$. Then we define
	\[\left|A\right|\coloneqq\sqrt{A\circ A^*}.\]
\end{notation}
\begin{definition}[compact]
	A bounded linear operator $A\colon V\to V$ is \textit{Hilbert--Schmidt} if and only if there is an orthonormal basis $\{v_i\}_i$ of $V$ such that
	\[\sum_i\norm{Av_i}_V^2<\infty.\]
	If $A$ is Hilbert--Schmidt, we define $\norm A_{\mathrm{HS}}\coloneqq\left(\sum_i\norm{Av_i}_V^2\right)^{1/2}$.
\end{definition}
\begin{definition}[trace class]
	A bounded linear operator $A\colon V\to V$ is \textit{trace class} if and only if there is an orthonormal basis $\{v_i\}_i$ of $V$ such that
	\[\sum_i(\left|A\right|v_i,v_i)<\infty.\]
	If $A$ has trace class, we define $\norm A_{\tr}\coloneqq\sum_i(\left|A\right|v_i,v_i)$ and $\tr A\coloneqq\sum_i(Av_i,v_i)$.
\end{definition}
\begin{remark}
	The norms and the trace are independent of the choice of basis, but this is tricky to show.
\end{remark}
\begin{example}
	Let $X$ and $Y$ be compact Hausdorff spaces equipped with Radon measures. Choose $K\in C(X\times Y)\cap L^2(X\times Y)$. Then the integral operator $T\colon L^2(X)\to L^2(Y)$ defined by
	\[Tf(y)\coloneqq\int_YK(x,y)f(x)\,dx\]
	is bounded and Hilbert--Schmidt. One sees this by comparing $T$ to certain ``Riemann sum operators,'' which are finite rank.
\end{example}
\begin{remark}
	Trace class implies Hilbert--Schmidt.
\end{remark}
\begin{theorem}[Spectral]
	Let $A\colon V\to V$ be a self-adjoint Hilbert--Schmidt operator on a Hilbert space $V$. Then there is a decomposition
	\[V=\ker A\oplus\bigoplus_{\lambda\Spec A}V[\lambda],\]
	where $\Spec A\subseteq\RR$ is a discrete set, except possibly with a limit point at zero, and $V[\lambda]\coloneqq\ker(A-\lambda1_V)$ is finite-dimensional.
\end{theorem}
\begin{proof}[Sketch]
	Inductively strip one eigenvalue off at a time.
\end{proof}
\begin{corollary}
	Let $(\pi,V)$ be a unitary representation of a locally compact Hausdorff group $G$. Suppose that $\{f_n\}$ is an approximate identity in $C_c(G)$ such that $\pi(f_n)$ is Hilbert--Schmidt for all $n$. Then $V$ is a Hilbert direct sum of its irreducible subrepresentations, each of which has finite multiplicity.
\end{corollary}
\begin{proof}[Sketch]
	This is fairly technical. To show that the multiplicities are finite, for any irreducible $\pi'$, choose $f\in C_c(G)$ with $\pi'(f)\ne0$ so that there is a nonzero eigenvalue $\lambda$ of $\pi'(f)$. Then choose some $v'\in\pi'$, and evaluation at $v'$ shows
	\[\dim\op{Hom}_G(\pi',\pi)\le\dim\ker(f-\lambda;V),\]
	which is finite by the Spectral theorem.

	To exhibit the decomposition, we use Zorn's lemma. By a Zorn's lemma argument with orthogonal complements, it is enough to show that any nonzero subrepresentation $W\subseteq V$ admits an irreducible subrepresentation. We may as well assume $W=V$. Let's describe how to construct the irreducible subrepresentation: choose $f\in C_c(G)$ with $\pi(f)\ne0$. Then choose a nonzero eigenvalue $\lambda$ of $\pi(f)$; note $V[\lambda]$ is finite-dimensional. Define $M$ to be a nonzero subspace minimal among the intersections
	\[\{V[\lambda]\cap V':V'\subseteq V\text{ is a subrepresentation}\}\setminus\{\emp\}.\]
	The subrepresentation generated by any nonzero vector in $M$ will work!
\end{proof}
\begin{proposition}
	Let $X$ be a locally compact second countable Hausdorff space equipped with a Radon measure. Let $\mc{HS}(X)$ denote the Hilbert space of Hilbert--Schmidt operators on $L^2(X)$. Then $\mc{HS}(X)\cong L^2(X\times X)$.
\end{proposition}
\begin{proof}[Sketch]
	The map $L^2(X\times X)\to\mc{HS}(X)$ is given by sending a kernel function to the corresponding integral operator. For the inverse map, fix a basis $\{\varphi_i\}_i$ of $L^2(X)$. Then send $A\in\mc{HS}(X)$ to the function
	\[K_A(x,y)\coloneqq\sum_iA\varphi_i(x)\overline{\varphi_i(x)}.\]
	In particular, one can show $A\mapsto K_A$ is an isometry.
\end{proof}
\begin{remark}
	The proof shows that the function
	\[\sum_iA\varphi_i(x)\overline{\varphi_i(x)}\]
	does not depend on the choice of basis.
\end{remark}
\begin{proposition}
	Let $X$ be a locally compact second countable Hausdorff space equipped with a Radon measure. Let $A\colon L^2(X)\to L^2(X)$ be an integral operator given by a continuous kernel $K\in L^2(X\times X)$. Then
	\[\tr A=\int_XK(x,x)\,dx.\]
\end{proposition}
\begin{proof}[Sketch]
	This is actually pretty tricky. Intuitively, on a basis $\{\varphi_i\}$ of $L^2(X)$, we find that $\tr A$ equals
	\[\sum_i\langle A\varphi_i,\varphi_i\rangle=\sum_i\int_XA\varphi_i(x)\overline{\varphi_i(x)}\,dx.\]
	Rearranging the sum with the integral completes the proof.
\end{proof}
\begin{remark}
	In the application, we will have $X=[G]$ and $A$ will usually be given by $Rf$. Then one can use the Dixmier--Malliavin theorem to reduce to the case where $f$ is a convolution $f_1*f_2$, so the trace is a product of $Rf_1$ and $Rf_2$. In this case, convergence is easier to show.
\end{remark}

\section{Nonarchimedean Representation Theory}
Throughout, $v$ is a finite place.

\subsection{The Hecke Algebra}
\begin{definition}[smooth]
	A representation $V$ of $G(F_v)$ is \textit{smooth} if and only if each $v\in V$ is stabilized by a compact open subgroup of $G(F_v)$.
\end{definition}
\begin{remark}
	Let $(\pi,V)$ be a smooth finite-dimensional representation of $G(F_v)$, where $G$ is split reductive. Then we claim that $\pi$ factors through $G^{\mathrm{ab}}(F_v)$. This follows from the Iwahori decomposition: $\ker\pi$ contains an open subgroup, so it contains $K_n$ for some $n>0$, so it contains a compact open subgroup of each root space $N_\alpha$. Because $\ker\pi$ is normal, we may conjugate this subgroup to cover $N_\alpha$. Lastly, the root spaces generate $G^{\mathrm{der}}$.
\end{remark}
\begin{definition}[admissible]
	A representation $V$ of $G(F_v)$ is \textit{admissible} if and only if it is smooth and
	\[\dim V^K<\infty\]
	for each compact open subgroup $K\subseteq G(F_v)$.
\end{definition}
\begin{definition}[Hecke algebra]
	Let $\mc H_G$ denote the \textit{Hecke algebra}, which is the algebra of smooth distributions $\mu\colon C_c^\infty(G(F_v))\to\CC$. For any compact open $K\subseteq G(F_v)$, define $\mc H_{G,K}$ as those distributions which are $K$-invariant on the left and right. The theory of the Haar measure allows us to identify $\mc H_G$ with $C_c^\infty(G(F_v))$ (uniquely up to scalar), which we will frequently do.
\end{definition}
\begin{remark}
	For any smooth representation $V$ of $G(F_v)$, note that $\mc H_G$ acts on $V$ by
	\[\mu\cdot f\coloneqq\int_{G(F_v)}f(x)\,d\mu(x).\]
\end{remark}
\begin{notation}
	For each compact open subgroup $K\subseteq G(F_v)$, define $e_K\in\mc H_G$ as
	\[\frac1{\op{vol}(K;dx)}\cdot1_K\,dx\in\mc H_G,\]
	where $dx$ is a choice of Haar measure on $G(F_v)$.
\end{notation}
\begin{remark}
	If $V$ is smooth, $e_K$ acts by projection $V\to V^K$.
\end{remark}
\begin{lemma}
	Fix a smooth representation $V$ of $G(F_v)$. Then $V$ is irreducible if and only if $V^K$ is zero or irreducible as a representation of $\mc H_{G,K}$ for each $K\subseteq G(F_v)$.
\end{lemma}
\begin{proof}[Sketch]
	In the reverse direction, note that a nonzero factorization of $V$ remains nonzero after taking invariants for small enough $K$ by smoothness. In the forward direction, note that if $W\subseteq V^K$ is a nonzero subrepresentation, then $\mc H_G\cdot W=V$, so $\mc H_{G,K}\cdot W=V^K$.
\end{proof}
\begin{lemma}
	Fix smooth irreducible representations $V_1$ and $V_2$ of $G(F_v)$. Choose $K\subseteq G(F_v)$ for which $V_1^K$ and $V_2^K$ are nonzero. If $V_1^K\cong V_2^K$, then $V_1\cong V_2$.
\end{lemma}
\begin{proof}[Sketch]
	Choose an isomorphism $I\colon V_1^K\to V_2^K$, and fix nonzero $v\in V_1^K$. It is enough to check that the extension
	\[I(f\cdot v_1)\coloneqq f\cdot I(v_1)\]
	is a well-defined map $V_1\to V_2$. Thus, we want to show $f\cdot v_1=0$ implies $f\cdot I(v_1)=0$. This is tricky: one can show $e_Kf_0f_1\cdot I(v_1)=0$ for any $f_0$, so $\mc H_{G,K}\cdot fI(v_1)=0$, but $\mc H_{G,K}\cdot fI(v_1)=\left(\mc H_G\cdot fI(v_1)\right)^K$, so $\mc H_G\cdot fI(v_1)\subsetneq V_2^K$, so $fI(v_1)=0$ by irreducibility!
\end{proof}

\subsection{Parabolic Induction and Jacquet Modules}
\begin{definition}[contragredient]
	Let $V$ be a smooth representation of $G(F)$. Then define $V^\lor$ to consist of the smooth linear functionals on $V$, equipped with the natural $G(F)$-action.
\end{definition}
\begin{remark}
	If $V$ is admissible, $V^\lor$ is admissible, and $V\cong V^{\lor\lor}$.
\end{remark}
\begin{definition}[parabolic induction]
	Let $P\subseteq G$ be a parabolic subgroup with Levi decomposition $P=MN$. For a representation $(\pi,V)$ of $M(F_v)$, we define $\op{Ind}_P^G\pi$ to be the representation
	\[\left\{\varphi\in C^\infty(G(F_v),V):\varphi(mng)=\delta_P^{1/2}(m)\cdot m\varphi(g)\right\},\]
	where $G(F_v)$ acts by right translation.
\end{definition}
\begin{remark}
	If $\pi$ is smooth, then $\op{Ind}_P^G\pi$ is smooth. If $\pi$ is admissible, then $\op{Ind}_P^G\pi$ is admissible.
\end{remark}
\begin{remark}
	If $\pi$ is unitary, then $\op{Ind}_P^G\pi$ is unitary. In short, choose a maximal compact subgroup $K\subseteq G(F_v)$ with $G(F_v)=P(F_v)K$. Then define the inner product
	\[\langle\varphi,\psi\rangle_{\op{Ind}_P^G\pi}\coloneqq\int_{K}\langle\varphi(x),\psi(x)\rangle\,dx.\]
	(The same inner product can be used to show that $\left(\op{Ind}_P^G\pi\right)^\lor\cong\op{Ind}_P^G\pi^\lor$.) To see that this is invariant, use the Iwasawa decomposition and rewrite the integral as $\int_{G(F_v)}\langle\varphi(x),\psi(x)\rangle f(x)\,dx$, where $f\in C_c^\infty(G(F_v))$ satisfies $\int_{P(F_v)}f(pg)\,dp=1$ for all $g$.
\end{remark}
\begin{definition}[Jacquet module]
	Let $P\subseteq G$ be a parabolic subgroup with Levi decomposition $P=MN$, and let $\psi\colon N(F_v)\to\CC^\times$ be a character. For a representation $V$ of $G(F_v)$, define the \textit{Jacquet module} $J_{\psi,N}(V)$ as
	\[V/V(N,\psi)\otimes\delta_P^{-1/2},\]
	where $V(N,\psi)\coloneqq\op{span}\{n\varphi-\psi(n)\varphi:\varphi\in V,n\in N(F_v)\}$. If $\psi=1$, we may omit it from the notation. We may omit $N$ from the notation. Note that $J_{\psi,N}(V)$ is a representation of $M(F)$.
\end{definition}
\begin{remark}
	If $V$ is smooth, then $J(V)$ is smooth.
\end{remark}
\begin{proposition}
	Let $P\subseteq G$ be a parabolic subgroup with Levi decomposition $P=MN$. The functors $J\colon\op{Rep}G(F_v)\to\op{Rep}M(F_v)$ and $\op{Ind}_P^G\colon\op{Rep}P(F_v)\to\op{Rep}G(F_v)$ are adjoint.
\end{proposition}
\begin{proof}[Sketch]
	Note that
	\[\op{Hom}_{G(F_v)}\left(V,\op{Ind}_P^GW\right)\cong\op{Hom}_{P(F_v)}(V,W)\cong\op{Hom}_{M(F_v)}(J(V),W).\]
	The first isomorphism holds by ``currying,'' sending $f\colon V\to\op{Ind}_P^GW$ to $v\mapsto f(v)(1)$. The second isomorphism holds because the action of $N(F_v)$ on $W$ is trivial.
\end{proof}
\begin{remark}
	The functor $\op{Ind}_P^G$ is exact on smooth representations. We already have left-exactness by the adjunction, so the only issue is right-exactness. The main idea is to write
	\[\op{Ind}_P^G(V)\cong\bigoplus_{s\in P(F_v)\backslash G(F_v)/K}V^{M(F_v)\cap sKs^{-1}},\]
	where $K\subseteq G(F_v)$ is some compact open subroup.
\end{remark}
\begin{lemma}
	Let $V$ be a smooth representation of $N(F_v)$. Then
	\[V(N)=\bigcup_{K\subseteq N(F_v)}\ker(R_K\colon V\to V),\]
	where the union is over compact open subgroups $K\subseteq N(F_v)$, and $R_K$ is the averaging operator.
\end{lemma}
\begin{proof}[Sketch]
	In one direction, note that $nv-v$ is in $\ker R_K$ for some $K$ by choosing $K\subseteq N(F_v)$ large enough to contain $n$. Conversely, if $\int_Knv\,dn=0$, then suppose that $v$ is fixed by some compact open $K'\subseteq K$, and then $v\simeq\sum_{s\in K/K'}(sv-v)$.
\end{proof}
\begin{remark}
	The functor $J$ is exact on smooth representations. We already have right-exactness. Left-exactness basically follows from the lemma, which explains why $-(N)$ is right-exact.
\end{remark}
\begin{lemma}[transitivity]
	Let $P\subseteq G$ be a parabolic with Levi decomposition $P=MN$, and let $P'\subseteq M$ be a parabolic with Levi decomposition $P'=M'N'$. Then
	\[\mathrm{Ind}_P^G\circ\mathrm{Ind}_{P'}^{M'}=\mathrm{Ind}_{P'N}^G.\]
\end{lemma}
\begin{proof}[Sketch]
	By the adjunction, it is enough to check a corresponding statement on Jacquet modules, which holds basically because the unipotent radical of $P'N$ is $N'N$.
\end{proof}
\begin{proposition}
	Let $P\subseteq G$ be a parabolic with Levi decomposition $P=MN$. If $V$ is an admissible representation of $G(F_v)$, then $J(V)$ is an admissible representation of $M(F_v)$.
\end{proposition}
\begin{proof}[Sketch]
	We use the Iwahori decomposition; set $M_n\coloneqq M(F_v)\cap K_n$ and $N_n^\pm\coloneqq N^\pm(F_v)\cap K_n$ for brevity. We will show that $J(V)^{M_n}$ is finite-dimensional for each $n$. By lifting finite-dimensional subspaces of $J(V)^{M_n}$ to $V$ and conjugating, we see that $J(V)^{M_n}$ can be covered by (translates of) $J(V)^{M_nN_n^-}$.

	Thus, it is enough to show that $V^{M_nN_n^-}$ and $V^{K_n}$ have the same image in $J(V)$. To see the claim, note that the difficult inclusion will follow from the claim
	\[R_{K_n}(v)\stackrel?\equiv v\pmod{\ker R_{N_n^+}}\]
	for all $v\in V^{M_nN_n^-}$; one shows this by using the Iwahori decomposition to compute $R_{K_n}(v)=R_{N_n^+}(v)$.
	%
	% To use the claim, we will show that $\dim J(V)^{M_n}$ is finite-dimensional for each $n$ by showing that $\dim U\le\dim V^{K_n}$ for each finite-dimensional subspace $U\subseteq J(V)^{M_n}$. Then lift $U$ to some $\widetilde U\subseteq V^{M_n}$ of the same dimension. By smoothness, $\widetilde U\subseteq V^{M_nN'}$ for some compact open $N'\subseteq N^-(F_v)$, but we can conjugate $N^-(F_v)$ to land in $N_n^-$.
\end{proof}

\subsection{Supercuspidal Representations}
\begin{definition}[supercuspidal, quasicuspidal]
	A representation $V$ of $G(F_v)$ is \textit{quasicuspidal} if and only if it is smooth, and its matrix coefficients have compact support modulo $Z(G(F_v))$. If it is further admissible, then $V$ is \textit{supercuspidal}.
\end{definition}
\begin{proposition}
	Let $V$ be an irreducible quasicuspidal representation of $G(F_v)$. Then $V$ is supercuspidal.
\end{proposition}
\begin{proof}[Sketch]
	We proceed by contraposition, supposing that $V$ fails to be admissible, and then we will exhibit a matrix coefficient which is not compactly supported modulo $Z(G(F_v))$. We are given a compact open subgroup $K\subseteq G(F_v)$ for which $V^K$ is infinite-dimensional; let $R_K\colon V\to V^K$ be the averaging operator.

	Let's just describe the construction of the matrix coefficient. Fix some nonzero $\varphi\in V$, and then there is an infinite set for which $S\subseteq G(F_v)$ such that $\{R_K(g\varphi):g\in S\}$ is a basis of $V^K$. Define $\varphi^\lor\in V^\lor$ by sending $\ker R_K$ to zero and sending $R_K(g\varphi)$ to $1$ for each $g\in S$. Then $x\mapsto\langle x\varphi,\varphi^\lor\rangle$ is the requested matrix coefficient.
\end{proof}
\begin{definition}[algebraically cuspidal]
	A representation $V$ of $G(F_v)$ is \textit{algebraically cuspidal} if and only if $J_N(V)=0$ for all $N$, as $N$ varies over the unipotent radicals of parabolic subgroups of $G$.
\end{definition}
\begin{proposition}
	Let $V$ be a smooth finitely generated representation of $G(F_v)$. Then there exists a para\-bolic subgroup $P\subseteq G$ with Levi decomposition $P=MN$ and some algebraically cuspidal smooth irreducible representation $W$ of $M(F_v)$ for which there is a nonzero intertwining map $V\to\op{Ind}_P^GW$.
\end{proposition}
\begin{proof}[Sketch]
	Let $P\subseteq G$ be a maximal (standard) parabolic admitting Levi decomposition $P=MN$ for which $J_N(V)\ne0$; if no such proper parabolic exists, then $V$ is already algebraically cuspidal, and there is nothing to do. Otherwise, by transitivity of the Jacquet module, $J_N(V)$ is algebraically cuspidal. Now, $J_N(V)$ is a finitely generated nontrivial representation of $M(F_v)$, so it admits a maximal subrepresentation by a Zorn's lemma argument, so it admits an irreducible quotient $W$. Note $W$ is algebraically cuspidal as well, so we conclude by Frobenius reciprocity.
\end{proof}
\begin{remark}
	If $V$ is actually irreducible, then the pair $(M,W)$ is unique (up to suitable conjugation).
\end{remark}
\begin{example}
	It's worthwhile to explicate this for $G=\op{GL}_2$. Let $V$ be a smooth irreducible representation. If $J(V)=0$, then $V$ is already supercuspidal. Otherwise, $J(V)$ is nonzero and thus admits a one-dimensional quotient $\chi$; the map $J(V)\to\chi$ induces a nonzero intertwining map $V\to\op{Ind}_B^G\chi$.
\end{example}
\begin{theorem}
	Let $V$ be a smooth representation of $G(F_v)$. Then $V$ is algebraically cuspidal if and only if $V$ is quasicuspidal.
\end{theorem}
\begin{proof}[Sketch]
	This essentially follows from the Cartan decomposition, so fix an Iwasawa decomposition $G(F_v)=P_0(F_v)K_0$ and maximal split torus $A_0\subseteq P_0$. Let $A\subseteq Z(G)$ be a maximal split torus.
	\begin{itemize}
		\item In the forward direction, we would like to show that a chosen matrix coefficient $g\mapsto\langle g\varphi,\varphi^\lor\rangle$ is compactly supported modulo $Z(G)$. By the Cartan decomposition, it is enough to choose some $\alpha\in\Delta$ and show
		\[\left\langle K_0\varpi^\lambda K_0\varphi,\varphi^\lor\right\rangle\stackrel?=0\]
		for $\lambda\in X_*(A_0)^+$ with $\alpha(\lambda)$ large.
		
		Let's describe the construction. In short, set $P\coloneqq P_{\Delta\setminus\{\alpha\}}$ with standard Levi decomposition $P=MN$, and then we may use that $V$ is algebraically cuspidal to choose an open compact subgroup $N_\varphi\subseteq N(F_v)$ for which $R_{N_\varphi}(\varphi)=0$. On the other hand, use smoothness to choose an open compact $N_{\varphi^\lor}\subseteq N(F_v)$ fixing every vector in $K_0\cdot\varphi^\lor$. It is now enough to choose $\lambda$ large enough so that $\varpi^\lambda N_\varphi\varpi^{-\lambda}\subseteq N_{\varphi^\lor}$.

		\item In the reverse direction, we would like to show that $V=V(N)$, where $N$ is the unipotent radical of the standard parabolic $P\coloneqq P_{\Delta\setminus\{\alpha\}}$. Well, choose $\varphi\in V$, and we want to show that $R_{N_\varphi}(\varphi)=0$ for some compact open $N_\varphi\subseteq N(F_v)$.

		The key idea is to test for $R_{N_\varphi}(\varphi)=0$ against given linear functionals $\varphi^\lor\in V^\lor$, which turns out to be enough by a smoothness argument. Let's explain where the quasicuspidal hypothesis is used: by reversing the argument of the previous point, we see that
		\[\left\langle\varpi^\lambda\varphi,K_0\varphi^\lor\right\rangle=0\]
		for $\alpha(\lambda)\gtrsim_{\varphi^\lor}1$. Then $N_{\varphi^\lor}\coloneqq\varpi^{-\lambda}(K_0\cap N(F_v))\varpi^\lambda$ has $\varphi^\lor\left(R_{N_{\varphi^\lor}}\varphi\right)=0$.
		\qedhere
	\end{itemize}
\end{proof}
\begin{corollary}
	Let $V$ be a smooth irreducible representation of $G(F_v)$. Then $V$ is admissible.
\end{corollary}
\begin{proof}[Sketch]
	Embed $V$ into the parabolic induction of an algebraically cuspidal irreducible representation. The latter is admissible.
\end{proof}
\begin{theorem}[Uniform admissibility]
	Fix a compact open subgroup $K$ of $G(F_v)$. Then there is a constant $D(K)$ such that
	\[\dim V^K\le D(K)\]
	for all smooth irreducible representations $V$ of $G(F_v)$.
\end{theorem}
\begin{proof}[Sketch]
	Omitted for time.\todo{}
\end{proof}

\subsection{Matrix Coefficients}
\begin{definition}[square integrable]
	Suppose $G$ is reductive. A smooth representation $V$ of $G(F_v)$ is \textit{square integrable} if and only if its matrix coefficients are in $L^2(G(F_v))$; it is \textit{essentially square integrable} if and only if its matrix coefficients are in $L^2\left(G(F_v)^1\right)$ upon restriction.
\end{definition}
\begin{definition}[tempered]
	Suppose $G$ is reductive. A smooth representation $V$ of $G(F_v)$ is \textit{tempered} if and only if its matrix coefficients are in $L^{2+\varepsilon}(G(F_v))$ for all $\varepsilon>0$; it is \textit{essentially square integrable} if and only if its matrix coefficients are in $L^{2+\varepsilon}\left(G(F_v)^1\right)$ for all $\varepsilon>0$ upon restriction.
\end{definition}
\begin{lemma}[Orthogonality]
	Suppose that $V$ and $W$ are admissible irreducible representations of $G(F_v)$ such that $V\not\cong W^\lor$ on $G(F_v)^1$, and suppose $V$ is supercuspidal. For $v\in V$ and $v^\lor\in V^\lor$ and $w\in W$ and $w^\lor\in W^\lor$, we have
	\[\int_{G(F_v)^1}\langle gv,v^\lor\rangle\langle gw,w^\lor\rangle\,dg=0.\]
\end{lemma}
\begin{proof}[Sketch]
	The assumption that $V$ is supercuspidal implies that the integrand is compactly supported. Fix $v^\lor$ and $w^\lor$. As a function in $v$ and $w$, this map is a $G(F_v)$-invariant pairing $V\times W\to\CC$, which must vanish.
\end{proof}
\begin{lemma}[Orthogonality]
	Let $V$ be a supercuspidal irreducible representation of $G(F_v)$. Then there is a constant $c_\pi>0$ for which
	\[\int_{G(F_v)^1}\langle gv,v^\lor\rangle\langle w,gw^\lor\rangle\,dg=c_\pi\langle v,w^\lor\rangle\langle w,v^\lor\rangle.\]
\end{lemma}
\begin{proof}[Sketch]
	Fixing $v^\lor$ and $w$ shows that this must be proportional to $\langle v,w^\lor\rangle$. Similarly, that constant of proportionality must be proporitional to $\langle v^\lor,w\rangle$. Positivity follows by plugging in some vectors.
\end{proof}
\begin{definition}[character]
	Let $(\pi,V)$ be an admissible representation of $G(F_v)$. For any $f\in C_c^\infty(G(F_v))$, choose a compact open subgroup $K\subseteq G(F_v)$ for which $f$ is bi-invariant under $K$; then choose any finite-dimensional subspace $W$ containing $V^K$ and define
	\[\tr\pi(f)\coloneqq\tr(\pi(f);W).\]
	This definition gives a well-defined \textit{character} distribution $\tr\pi\colon C_c^\infty(G(F_v))\to\CC$.
\end{definition}
\begin{theorem}[Harish-Chandra]
	Suppose $\op{char}F=0$. Let $G^{\mathrm{rs}}\subseteq G$ be the Zariski open subset of regular semisimple elements. Then there is $\Theta_\pi\in C^\infty(G^{\mathrm{rs}}(F_v))$ such that
	\[\tr\pi(f)=\int_{G(F_v)}\Theta_\pi(x)f(x)\,dx\]
	for all $f\in C_c^\infty(G(F_v))$.
\end{theorem}
\begin{proof}[Sketch]
	Hard, so entirely omitted.
\end{proof}
\begin{proposition}
	Let $\{(\pi_i,V_i)\}_i$ be a finite collection of pairwise-non-isomorphic irreducible admissible represenations of $G(F_v)$. Then the distributions $\{\tr\pi_i\}_i$ are linearly independent.
\end{proposition}
\begin{proof}[Sketch]
	We pass to a finite-dimensional statement. Choose compact open $K\subseteq G(F_v)$ small enough so that $V_i^K\ne0$ for each $i$. Now, let $A$ be the image of the action map
	\[\mc H_{G,K}\to\prod_i\op{End}_\CC\left(V_i^K\right).\]
	The finite-dimensional algebra $A$ admits the $\{V_i\}_i$ as a finite set of irreducible modules for which $A$ acts faithfully on the set. Thus, $A$ is semisimple, and the proposition follows from orthogonality of characters for the irreducible modules of $A$.
\end{proof}
\begin{corollary}
	Suppose that $G$ is split reductive, and let $\iota\colon G\to G$ be the Cartan involution. Let $(\pi,V)$ be an irreducible admissible representation of $G(F_v)$. Then $V^\lor$ is isomorphic to the representation $g\mapsto\pi(\iota g)$.
\end{corollary}
\begin{proof}[Sketch]
	It is enough to compare the characters. Let $\widetilde V$ be the latter representation. Fix $f\in C_c^\infty(G(F_v))$. On one hand, $\Theta_{V^\lor}(f)=\Theta_V\left(f\circ(-)^{-1}\right)$ by rearranging the pairing. On the other hand, $\Theta_{\widetilde V}(f)=\Theta_V(f\circ\iota)$. From here, one can check (on the Lie algebra) that $g$ and $\iota(g)^{-1}$ are conjugate for any $g\in G(F_v)$.
\end{proof}
\begin{proposition}
	Let $(\pi,V)$ be an irreducible supercuspidal representation of $G(F_v)$. Suppose $Z_G(F_v)$ is compact. For any $f\in C_c^\infty(G(F_v))$, there is $f_\pi\in C_c^\infty(G(F_v))$ such that
	\[\pi'(f_\pi)=\begin{cases}
		\pi(f) & \text{if }\pi'\cong\pi, \\
		0 & \text{otherwise},
	\end{cases}\]
	for any irreducible admissible representation $\pi'$ of $G(F_v)$.
\end{proposition}
\begin{proof}[Sketch]
	The argument present in Getz--Hahn has a gap.
\end{proof}
\begin{proposition}
	Let $(\pi,V)$ be an irreducible supercuspidal representation of $G(F_v)$. Then there is $f_\pi\in C_c^\infty(G(F_v))$ with the following property: for any irreducible $(V',\pi')$, if $\tr\pi'(f_\pi)\ne0$, there is a character $\chi$ on $G(F_v)$ for which $\pi'\cong\pi\otimes\chi$.
\end{proposition}
\begin{proof}[Sketch]
	Let $f_\pi$ be a matrix coefficient of $V$. Orthogonality implies that $\tr\pi'(f_\pi)=0$ unless $V$ and $V'$ are isomorphic as representations of $G(F_v)^1$. In this latter case, identify $V$ and $V'$ as vector spaces, and then one can show that $\chi(g)\coloneqq\pi'(g)\pi(g)^{-1}$ factors through $G(F_v)\to G^{\mathrm{ab}}(F_v)$, which is covered by the center, so $\chi$ is the desired character.
\end{proof}
\begin{remark}
	Characters of the principal series live in a continuous family, so one cannot single out the principal series with a finite collection of functions. One shows the claimed continuity by a direct calculation relating to the character of an induction to the original representation.\todo{}
\end{remark}

\subsection{The Whittaker Model}
\begin{definition}[generic]
	Let $G$ be a quasi-split reductive group over $F_v$. Let $B$ be a Borel subgroup of $G$ with unipotent radical $N$, and let $T\subseteq B$ be a maximal split torus. A character $\psi\colon N(F_v)\to\CC^\times$ is \textit{generic} if and only if $\psi|_{N_\alpha(F)}\ne1$ for each $\alpha\in\Delta$.
\end{definition}
\begin{definition}[Whittaker functional]
	Let $G$ be a quasi-split reductive group over $F_v$, and let $B$ be a Borel subgroup with unipotent radical $N$. Choose $\psi\colon N(F_v)\to\CC^\times$.
	\begin{itemize}
		\item If $v$ is finite, a \textit{$\psi$-Whittaker functional} is an element of $\op{Hom}_{N(F_v)}(V,\psi)$.
		\item If $v$ is infinite, a $\psi$-Whittaker functional is a continuous linear functional in $\op{Hom}_{N(F_v)}(V_{\mathrm{sm}},\psi)$, where $V_{\mathrm{sm}}$ is given its natural Fr\'echet topology.
	\end{itemize}
\end{definition}
\begin{remark}
	In the finite case, a $\psi$-Whittaker functional exists if and only if $J_\psi(V)$ admits a functional.
\end{remark}
\begin{definition}[generic]
	Let $G$ be a quasi-split reductive group over $F_v$. A smooth irreducible representation $V$ is \textit{generic} if and only if it admits a $\psi$-Whittaker functional for some generic $\psi$.
\end{definition}
\begin{theorem}
	Let $G$ be a quasi-split reductive group over $F_v$, and let $B$ be a Borel subgroup with unipotent radical $N$. If $\psi\colon N(F_v)\to\CC^\times$ is generic, the space of $\psi$-Whittaker functionals of any smooth irreducible representation of $V$ is at most one.
\end{theorem}
\begin{proof}[Sketch]
	We only sketch this when $G$ is split and $v$ is finite. Let $\mc H'\subseteq\mc H_G$ be the eigenvectors of $N(F)\times N(F)$ with eigenvalue $\psi_N\boxtimes\psi_N^{-1}$. By a computation with the Bruhat decomposition, one finds that the Cartan involution $\iota$ acts trivially on $\mc H'$; this uses the hypothesis that $\psi$ is generic.

	Now, choose two nonzero $\psi$-Whittaker functionals $\lambda_1,\lambda_2\colon V\to\CC$. It is enough to show that $\ker\lambda_1=\ker\lambda_2$. By identifying $V$ with $V^\lor$ (as vector spaces with different $G(F_v)$-actions), we see that $\lambda_i$ becomes a $\psi^{-1}$-Whittaker functional on $V^\lor$. Now, the distributions
	\[\varphi\mapsto\langle\lambda_1*\varphi,\lambda_1'\rangle\qquad\text{and}\qquad\varphi\mapsto\langle\lambda_2*\varphi,\lambda_2'\rangle\]
	both live in $\mc H'$, and a direct calculation shows that they differ by the action of $\iota$, so they must actually be equal! This can be rearranged into the desired result.\todo{}
\end{proof}
\begin{definition}[Whittaker function]
	Let $G$ be a quasi-split reductive group over $F_v$, and let $B$ be a Borel subgroup with unipotent radical $N$. For any generic $\psi\colon N(F_v)\to\CC^\times$, a \textit{$\psi$-Whittaker function} is an element of the space $\mc W(\psi)\coloneqq C^\infty((N(F_v),\psi)\backslash G(F_v))$.
\end{definition}
\begin{remark}
	By Frobenius reciprocity, the presence of a $\psi$-Whittaker functional $\lambda$ of $V$ is equivalent to admitting a Whittaker model $V_{\mathrm{sm}}\to\mc W(\psi)$ (given by sending $v$ to $W_v\colon g\mapsto\lambda(gv)$).
\end{remark}
\begin{theorem}
	Every supercuspidal representation of $\op{GL}_n(F_v)$ is generic.
\end{theorem}
\begin{proof}[Sketch]
	We will show this for $n=2$ shortly.
\end{proof}
\begin{proposition}
	Any non-generic smooth representation $V$ of $\op{GL}_2(F_v)$ factors through $\det$.
\end{proposition}
\begin{proof}[Sketch]
	This is pretty tricky. Consider the action $\widehat\pi$ of $C_c^\infty(F)$ on $V$ by
	\[\widehat\pi(f)v\coloneqq\int_F\widehat f(x)\cdot\begin{bmatrix}
		1 & -x \\ & 1
	\end{bmatrix}v\,dx,\]
	where the Fourier transform is taken with respect to some fixed $\psi\colon F_v\to\CC^\times$. This action is (co)smooth, and $1_{a+\mf p^k}$ acts by $\psi_a1_{\mf p^?}$ (up to scalar). Now, define a sheaf $\mc F_V$ on $F$ on the base of compact open subsets given by $U\mapsto1_UV$. A direct calculation shows that the stalk of $\mc F_V$ at $a\in F$ is $J_{\psi_a}(V)$, so $\mc F_V$ must be a skyscraper sheaf! Thus, $V=J(V)$, so $N(F_v)$ acts trivially on $V$. Taking conjugates shows that $\op{SL}_2(F_v)$ acts triivally on $V$.
\end{proof}

\section{Representation Theory of \texorpdfstring{$\op{GL}_2$}{GL(2)}}

\subsection{The Principal Series}
\begin{lemma}
	Let $P$ be a parabolic subgroup of $G$ with Levi decomposition $P=MN$. For any character $\chi\colon M(F_v)\to\CC^\times$, the map $R\colon C_c^\infty(G(F_v))\to\op{Ind}_P^G\chi$ defined by
	\[R_\varphi(x)\coloneqq\int_{P(F_v)}\varphi\left(p^{-1}x\right)\delta^{1/2}\chi(x)\,dx\]
	is surjective.
\end{lemma}
\begin{proof}[Sketch]
	Fix an Iwasawa decomposition $G(F_v)=P(F_v)K$. Given $f\in\op{Ind}_P^G\chi$, define $\varphi\colon G(F_v)\to\CC$ by $\varphi\coloneqq f1_K$. Then $R_\varphi\simeq f$.
\end{proof}
\begin{lemma}
	Fix a character $\xi\colon G(F_v)\to\CC^\times$. Any eigen-distribution $D\colon C_c^\infty(G(F_v))\to\CC$ satisfying $D(L_{g^{-1}}f)=\overline\xi(g)D(f)$ has is proportional to the distribution
	\[f\mapsto\int_{G(F_v)}\xi(h)f(h)\,d_\ell h.\]
\end{lemma}
\begin{proof}[Sketch]
	By twisting $D$, we may assume $\xi=1$. Now, fix a compact open subgroup $K\subseteq G(F_v)$, and normalize $d_\ell h$ so that $K$ has volume one. Then the uniqueness of the left-invariant Haar measure implies that
	\[Tf=T(1_K)\int_{G(F_v)}f(h)\,d_\ell h,\]
	as required.
\end{proof}
\begin{proposition}
	Let $G$ be a split reductive group with Borel $B=TN$. For each character $\chi\colon T(F_v)\to\CC^\times$ and generic $\psi\colon N(F_v)\to\CC^\times$,
	\[\dim J_\psi\left(\op{Ind}_B^G\chi\right)=1.\]
\end{proposition}
\begin{proof}[Sketch]
	This is basically Mackey theory. To start, note that
	\[f\mapsto\int_{N(F_v)}f(wn)\psi(n)^{-1}\,dn\]
	provides the required functional. (Even if $f$ does not have compact support, one can regularize the integral because $f$ will be constant ``at $\infty$,'' so the integral against $\psi^{-1}$ vanishes.) It remains to upper-bound. Via the projection $C_c^\infty(G(F_v))\to\op{Ind}_B^G\chi$, it is enough to show that
	\[\dim\left\{D\in C_c^\infty(G(F_v))^\lor:L_bD=\delta^{1/2}\chi(b)D\text{ and }R_nD=\psi(n)^{-1}D\right\}\le1.\]
	To avoid a complicated induction on the Bruhat decomposition, we only show this for $G=\mathrm{GL}_2$. Then the Bruhat decomposition simply gives
	\[0\to C_c^\infty(B(F_v))^\lor\to C_c^\infty(\op{GL}_2(F_v))^\lor\to C_c^\infty(B(F_v)wN(F_v))^\lor\to0,\]
	where $w$ is the long Weyl element. Now, any suitable eigen-distribution $D$ is uniquely determined up to scalar when restricted to $B(F_v)wN(F_v)\cong B(F_v)\times N(F_v)$. Thus, any two suitable eigen-distributions are equal up to an eigen-distribution from $B(F_v)$; but there are no suitable eigendistributions on $B(F_v)$, which we see by checking that the previous lemma does not provide a suitable eigen-distribution if we just use one of the eigen-conditions.
\end{proof}
\begin{remark}
	The same argument shows that $\dim J\big(\op{Ind}_B^G\chi\big)\le\op{rank}G$, and I think equality holds. For $G=\op{GL}_2$, one can show by hand that the functionals $f\mapsto f(1)$ and
	\[f\mapsto\int_{F_v}\left(f\left(w\begin{bmatrix}
		1 & x \\ & 1
	\end{bmatrix}\right)-\left|x\right|^{-1}\chi_1^{-1}\chi_2(x)1_{x\notin\OO_v}\varphi(1)\right)\,dx\]
	are both well-defined and linearly independent.
\end{remark}
\begin{proposition}
	Let $G$ be a split reductive group with Borel $B=TN$. If $\chi\colon T(F_v)\to\CC^\times$ has trivial stabilizer via the $(W,\cdot)$-action, then $\op{Ind}^G_B\chi$ is irreducible.
\end{proposition}
\begin{proof}[Sketch]
	We only show this for $G=\mathrm{GL}_2$. Suppose that $V\coloneqq\op{Ind}_B^G\chi$ lives in a short exact sequence
	\[0\to V'\to V\to V''\to0.\]
	For dimension reasons, $J_\psi(V')=0$ or $J_\psi(V'')=0$. By passing to the contragredient, we may take $J_\psi(V')=0$. Now, $V'$ factors through $\det$, so $V'$ admits an eigenvector, so $V$ admits a one-dimensional subspace; let $f$ be an eigenvector. Then $f$ satisfies
	\[f(bg)=\delta^{1/2}\chi(b)\rho(\det g)f(1)\]
	for some character $\rho\colon F^\times\to\CC^\times$. Plugging in $b=g^{-1}=\op{diag}(a,1/a)$ reveals that $\chi_1\chi_2^{-1}=\left|\cdot\right|^{-1}$, so $\chi$ is fixed by the $(W,\cdot)$-action.
\end{proof}
\begin{remark}
	Even in the exceptional cases $\chi_1\chi_2^{-1}=\left|\cdot\right|^{-1}$, the same argument shows that $\op{Ind}_B^G\chi$ contains an eigenvector with irreducible quotient. Dually, when $\chi_1\chi_2^{-1}=\left|\cdot\right|$, we see $\op{Ind}_B^G\chi$ contains a one-dimensional quotient with irreducible kernel.
\end{remark}
\begin{proposition}
	Let $G$ be a split reductive group with Borel $B=TN$. If $\chi,\chi'\colon T(F_v)\to\CC^\times$ have $\op{Ind}_B^G\chi\cong\op{Ind}_B^G\chi'$, then $\chi'\in W\chi$.
\end{proposition}
\begin{proof}[Sketch]
	This is Mackey theory.  Again, to avoid complicated inductions, we only show this for $G=\op{GL}_2$. Fix a nonzero morphism $\op{Ind}_B^G\chi\to\op{Ind}_B^G\mu$. By Frobenius reciprocity, we get a nonzero $B(F_v)$-invariant map $\op{Ind}_B^G\chi\to\delta^{1/2}\mu$. Via the projection $C_c^\infty(G(F_v))\to\op{Ind}_B^G\chi$, we receive a distribution $D$ satisfying 
	\[\begin{cases}
		L_bD=\delta^{-1/2}\chi(b)D, \\
		R_bD=\delta^{-1/2}\mu(b)^{-1}D,
	\end{cases}\]
	for $b\in B(F_v)$. If $D$ is nonzero on $B(F_v)wN(F_v)\cong B(F_v)\times N(F_v)$, then the eigenector conditions on $B(F_v)\times N(F_v)$ give $D$ an explicit formula, and then the remaining eigenvector condition on $T(F_v)$ forces $\mu={^w\chi}$. Otherwise $D$ is supported on $B(F_v)$, and one can compare the two eigenvector conditions to get $\mu=\chi$.
\end{proof}
\begin{proposition}
	Let $G$ be a split reductive group with Borel $B=TN$, and choose $\chi\colon T(F_v)\to\CC^\times$. If $\chi$ has trivial stabilizer via the $W$-action, then
	\[J\left(\op{Ind}_B^G\chi\right)\cong\delta^{1/2}\otimes\bigoplus_{w\in W}{^w\chi}.\]
\end{proposition}
\begin{proof}[Sketch]
	For a character $\mu\colon T(F_v)\to\CC^\times$, Frobenius reciprocity gives
	\[\op{Hom}_{T(F_v)}\left(J\left(\op{Ind}_B^G\chi\right),\delta^{1/2}\mu\right)\cong\op{Hom}_{G(F_v)}\left(\op{Ind}_B^G\chi,\op{Ind}_B^G\mu\right).\]
	By the previous proposition, this is nonzero exactly for $\mu\in W\chi$, so we have a quotient isomorphic to $\delta^{1/2}\otimes\bigoplus_{w\in W}{^w\chi}$. This is an isomorphism for dimension reasons.
\end{proof}
\begin{example}
	For $G=\mathrm{GL}_2$, let's describe what happens in the case $\chi={^w\chi}$. Then $J\big(\op{Ind}_B^G\chi\big)$ admits a unique quotient map to $\delta^{1/2}\chi$ (up to scalar); let the kernel be $\delta^{1/2}\mu$, so there is a short exact sequence
	\[0\to\delta^{1/2}\mu\to J\big(\op{Ind}_B^G\chi\big)\to\delta^{1/2}\chi\to0.\]
	If $\chi\ne\mu$, this splits, which is a contradiction. Instead, we must have $\mu=\chi$, and the short exact sequence must be non-split. Some explicit calculations on bases shows that we can give $J\big(\op{Ind}_B^G\chi\big)$ a basis so that it is isomorphic to
	\[t\mapsto\delta^{1/2}\chi(t)\begin{bmatrix}
		1 & v(t_1/t_2) \\ & 1
	\end{bmatrix}.\]
\end{example}

\subsection{Unramified Representations}
\begin{theorem}
	Let $G$ be a split reductive group, and let $K\subseteq G(F_v)$ be a maximal compact subgroup. Then $\mc H_K$ is commutative.
\end{theorem}
\begin{proof}[Sketch]
	Use Gelfand's trick. The relevant anti-involution is given by the (inverse of the) Cartan involution. For example, if $G=\mathrm{GL}_n$, then the anti-involution is $(-)^{-\intercal}$.
\end{proof}
\begin{definition}
	Fix a split reductive group $G$ with maximal compact subgroup $K\subseteq G(F_v)$. A smooth representation $V$ is \textit{spherical} if and only if $V^K\ne0$.
\end{definition}
\begin{remark}
	If $V$ is spherical and irreducible, then $V^K$ is an irreducible $\mc H_K$-module, so $\dim V^K=1$.
\end{remark}
\begin{theorem}[Satake]
	Let $G$ be a split reductive group with Borel $B=TN$ and Iwasawa decomposition $G(F_v)=B(F_v)K$. Let $\widehat T$ be the Langlands dual of $T$. Then
	\[\mc H_K\cong\OO(\widehat T)^W.\]
\end{theorem}
\begin{proof}[Sketch]
	We only argue this for $G=\mathrm{GL}_2$ with $K=\op{GL}_2(\OO_v)$, and we will argue some of this now. For each $k\ge0$, let $T_k\in\mc H_K$ be the indicator of $g\in M_2(\OO)$ with $(\det g)\OO=\mf p^k$; set $T\coloneqq T_1$ for brevity. Also, define $R$ to indiate $\varpi K$. An induction shows that $\mc H_K$ is spanned by the vectors $R^aT_k$ where $a\in\ZZ$ and $k\ge0$. By examining actions on $\PP^1_{\FF_q}$, one can show that
	\[K\op{diag}(\varpi,1)K=\begin{bmatrix}
		\varpi \\ & 1
	\end{bmatrix}K\sqcup\bigsqcup_{b\in\FF_q}\begin{bmatrix}
		\varpi & b \\ & 1
	\end{bmatrix}K.\]
	From here, some explicit calculation shows that $T_1T_k=T_{k+1}+qRT_{k-1}$ for $k\ge1$. It follows that the natural map $\CC[T,R,1/R]\to\mc H_K$ is a surjection, and this surjection will soon be checked to be an isomorphism. The source is isomorphic to $\CC[t_1+t_2,t_1t_2,1/t_1t_2]$, which we call good enough.
\end{proof}
\begin{example}
	Let $\alpha\colon F_v^\times\to\CC^\times$ be an unramified character. Then $\alpha$ is spherical. For $G=\op{GL}_2$, a direct calculation shows that $T_1$ acts by $(q+1)\alpha(\varpi)$ and $R$ acts by $\alpha(\varpi)^2$.
\end{example}
\begin{example}
	Let $\chi\colon T(F_v)\to\CC^\times$ be an unramified character. By the Iwasawa decomposition, $\op{Ind}_B^G\chi$ is unramified with the unique unramified vector $f\colon G(F_v)\to\CC$ given by $f(bk)=\delta^{1/2}\chi(b)$. For $G=\op{GL}_2$, a direct calculation shows that $T_1$ acts by $q^{1/2}(\chi_1(\varpi)+\chi_2(\varpi))$ and $R$ acts by $\chi_1(\varpi)\chi_2(\varpi)$.
\end{example}
\begin{theorem}
	Let $G$ be a split reductive group with Borel $B=TN$. Then every irreducible spherical representation $V$ is a subquotient of $\op{Ind}_B^G\chi$, where $\chi\colon T(F_v)\to\CC^\times$ is an unramified character (unique up to Weyl orbit from $V$).
\end{theorem}
\begin{proof}[Sketch]
	The fact that $\op{Ind}_B^G\chi$ admits a unique irreduible spherical subquotient follows from the fact that $\dim\big(\op{Ind}_B^G\chi\big)^K=1$. Now, the irreducible spherical representations embed into the characters of $\mc H_K$, so it merely remains to check that subquotients of $\op{Ind}_B^G\chi$ exhausts all such characters; this is an application of the Satake isomorphism.

	Let's explicate this for $G=\mathrm{GL}_2$. Any character $\xi\colon\mc H_K\to\CC$ is uniquely determined by $\xi(T)$ and $\xi(R)$, and we must have $\xi(R)\in\CC^\times$. Let $\alpha_1$ and $\alpha_2$ be the roots of the polynomial
	\[X^2-q^{-1/2}\xi(T)+\xi(R)=0.\]
	Then define $\chi_i\colon F^\times\to\CC^\times$ to be the unique unramified character for which $\chi_i(\varpi)=\alpha_i$. Then the character of $\mc H_K$ induced by $\op{Ind}_B^G\chi$ agrees with the character $\xi$. The only problem occurs when $\op{Ind}_B^G\chi$ fails to be irreducible; up to symmetry, this means $\chi_1\chi_2^{-1}=\left|\cdot\right|^{-1}$. Then $\xi$ is induced by the unramified character $\alpha\colon F_v^\times\to\CC^\times$ defined by $\alpha(\varpi)=q^{1/2}\alpha_2$.
\end{proof}
\begin{remark}
	This argument completes the proof of Satake's isomorphism for $G=\op{GL}_2$: indeed, it shows that the ring surjection $\CC[T,R,1/R]\to\mc H_K$ induces a surjective closed embedding $\Spec\mc H_K(\CC)\to\Spec\CC[T,R,1/R](\CC)$ with reduced target.
\end{remark}
\begin{theorem}[Cassleman--Shalika]
	Let $G$ be a split reductive group with Borel $B=TN$. Choose $\lambda\in X^*(T)_\CC$, and define $J(\lambda)$ to be the unique spherical irreducible quotient of
	\[\op{Ind}_B^Ge^{\langle H_T(-),\lambda\rangle}.\]
	Choose a generic character $\psi$, and suppose $W$ is a spherical vector in the Whittaker model of $J(\lambda)$.
	\begin{listalph}
		\item If $W\ne0$, then $W(1)\ne0$.
		\item If $\mu\in X^*(T)\setminus X^*(T)_+$, then $W(\varpi^\mu)=0$.
		\item If $\mu\in X^*(T)_+$, then
		\[W(\varpi^\mu)=\delta_B^{1/2}(\varpi^\mu)\chi_\mu\left(q^{-\lambda}\right)W(1).\]
	\end{listalph}
\end{theorem}
\begin{proof}[Sketch]
	We only argue this for $G=\op{GL}_2$, for which we give an ad-hoc argument. Let the representation be $\op{Ind}_B^G\chi$. We are tasked with computing the numbers $w_n\coloneqq W(\op{diag}(\varpi^n,1))$; the Iwasawa decomposition shows that $W$ is uniquely determined by these values. Set $\beta_i\coloneqq q^{-1/2}\chi_i(\varpi)$ for brevity. Using $W\left(\begin{bsmallmatrix}
		1 & x \\ & 1
	\end{bsmallmatrix}g\right)=\psi(x)W(g)$, one can show that $w_n=0$ if $n<0$. We know that $T\cdot W=q(\beta_1+\beta_2)W$. Comparing this with a direct expansion of $T\cdot W$ givess the recursion
	\[w_{n+2}=(\beta_1+\beta_2)w_{n+1}-\beta_1\beta_2w_n.\]
	Solving this recursion gives $w_n=\left(\beta_1^{n+1}-\beta_2^{n+1}\right)/(\beta_1-\beta_2)\cdot w_0$ and hence the result.
\end{proof}

\subsection{The Kirillov Model}
\begin{definition}[Kirillov model]
	Fix a generic smooth representation $V$ of $\op{GL}_2(F_v)$, and let $W_\bullet\colon V\to C^\infty((N(F_v),\psi)\backslash\op{GL}_2(F_v))$ be a Whittaker model. Then the \textit{Kirillov model} consists of the subspace of $C_c^\infty(F^\times)$ of functions
	\[\varphi_v(x)\coloneqq W_v(\op{diag}(x,1)).\]
\end{definition}
\begin{remark}
	For $\varphi$ in the Kirillov model, one has
	\[\begin{bmatrix}
		1 & b \\ & 1
	\end{bmatrix}\varphi(x)=\psi(bx)\varphi(x)\qquad\text{and}\qquad\begin{bmatrix}
		a \\ & d
	\end{bmatrix}\varphi(x)=\omega(d)\varphi(ax/d),\]
	where $\omega$ is the central character.
\end{remark}
\begin{lemma}
	Fix a generic smooth representation $V$ of $\op{GL}_2(F_v)$, which we identify with its Kirillov model. For any $\varphi\in V$, we have $\varphi(y)=0$ if $\left|y\right|\gtrsim_\varphi1$.
\end{lemma}
\begin{proof}[Sketch]
	By smoothness, $\varphi$ is stabilized by some subgroup $\left\{\begin{bsmallmatrix}
		1 & x \\ & 1
	\end{bsmallmatrix}:x\in\mf p^n\right\}$. Thus, for $x\in\mf p^n$, we see $(1-\psi(xy))\varphi(y)=0$, so $\varphi(y)=0$ as soon as $y\mf p^n\not\subseteq\ker\psi$.
\end{proof}
\begin{lemma}
	Fix a generic irreducible smooth representation $V$ of $\op{GL}_2(F_v)$, which we identify with its Kirillov model. The kernel of $V\to J(V)$ is exactly $C_c^\infty(F_v^\times)$.
\end{lemma}
\begin{proof}[Sketch]
	For one inclusion, recall $V(N)$ is spanned by vectors of the form $\begin{bsmallmatrix}
		1 & x \\ & 1
	\end{bsmallmatrix}\varphi-\varphi$, so it is enough to show that such functions vanish on small $y\in F_v^\times$. This follows because $\psi(xy)=1$ for small $y$.

	For the other inclusion, because the kernrel is nonzero, it is enough to show that $C_c^\infty(F_v^\times)$ is an irreducible as a $B(F_v)$-module. Let $V\subseteq C_c^\infty(F_v^\times)$ be a nonzero subspace, and we will check that $V$ contains sufficiently small indicators of any given $a\in F_v^\times$. By translation, we may assume there is $\varphi\in V$ with $\varphi(a)\ne0$. Then a direct expansion shows that some $f\in C_c^\infty(N(F_v))$ has
	\[f*\varphi=\widehat f\varphi,\]
	so just choose $\widehat f$ to indicate a small open neighborhood of $a$.
\end{proof}
\begin{proposition}
	Fix a generic irreducible smooth representation $V$ of $\op{GL}_2(F_v)$, which we identify with its Kirillov model.
	\begin{listalph}
		\item If $V=\op{Ind}_B^G\chi$ is an irreducible principal series with $\chi_1\ne\chi_2$, then $V$ contains $\left|\cdot\right|^{1/2}\chi_11_\OO$ and $\left|\cdot\right|^{1/2}\chi_21_\OO$.
		\item If $V=\op{Ind}_B^G\chi$ is an irreducible principal series with $\chi_1=\chi_2$, then $V$ contains $\left|\cdot\right|^{1/2}\chi_11_\OO$ and $v(t)\left|\cdot\right|^{1/2}\chi_11_\OO$.
		\item If $V\subsetneq\op{Ind}_B^G\chi$ is the codimension-one subspace when $\chi_1\chi_2^{-1}=\left|\cdot\right|$, then $V$ contains $\left|\cdot\right|^{1/2}\chi_11_\OO$.
	\end{listalph}
	Along with $C_c^\infty(F_v^\times)$, these vectors span $V$.
\end{proposition}
\begin{proof}[Sketch]
	We will only sketch (a) because the others are similar. Note $V(N)=C_c^\infty(F_v^\times)$, so it only remains to lift $J(V)$, which we recall is the representation $\delta^{1/2}\chi_1\oplus\delta^{1/2}\chi_2$ of $T(F_v)$.
	
	Choose $\varphi\in V$ such that $t\overline\varphi=\delta^{1/2}\chi(t)\overline\varphi$ in $J(V)$. We will show that $\varphi\simeq\left|\cdot\right|^{1/2}\chi_1$ for small $t$, which completes the argument. This is an intricate compactness argument. Well, given $t\in F^\times$, the lemma implies that
	\[\varphi(tu)=\delta^{1/2}\chi(\op{diag}(t,1))\varphi(u)\]
	holds for small $u$; say $\left|u\right|\le\varepsilon(t)$. By smoothness, $\varepsilon(t)$ can taken to be locally constant, so we may define $\varepsilon\coloneqq\min\{\varepsilon(t):t\in\varpi\OO^\times\}$.
\end{proof}

\subsection{Local Rankin--Selberg for \texorpdfstring{$\op{GL}_2$}{GL(2)}}
\begin{definition}[local zeta integral]
	For $\varphi\in C^\infty(F^\times)$, define
	\[Z(s,\varphi)\coloneqq\int_{F^\times}\varphi(y)\left|y\right|^{s-1/2}\,d^\times y,\]
	provided that the integral converges. For a character $\xi\colon F^\times\to\CC^\times$, define $Z(s,\varphi,\xi)\coloneqq Z(s,\varphi\xi)$.
\end{definition}
\begin{remark}
	We claim that $\{Z(s,\varphi):\varphi\in C_c^\infty(F_v^\times)\}=\CC\left[q^s,q^{-s}\right]$. The main point is that $Z(s,1_{1+\mf p^n})$ is a (positive) constant if $n>0$, and $Z(s,\op{diag}(a,1)\varphi)=\left|a\right|^{1/2}\left|a\right|^{-s}Z(s,\varphi)$ for any $a\in F^\times$.
\end{remark}
\begin{definition}[local $L$-factor]
	Let $V$ be a smooth generic irreducible representation of $\op{GL}_2(F_v)$.
	\begin{itemize}
		\item If $V$ is supercuspidal, define $L(s,\pi)\coloneqq1$.
		\item If $V=\op{Ind}_B^G\chi$ is an irreducible principal series, define $\alpha_i$ to be $\chi_i(\varpi)$ if $\chi_i$ is unramified and zero otherwise; then define $L(s,\pi)\coloneqq(1-\alpha_1q^{-s})^{-1}(1-\alpha_2q^{-s})^{-1}$.
		\item If $V\subsetneq\op{Ind}_B^G\chi$ is the codimension-one subspace when $\chi_1\chi_2^{-1}=\left|\cdot\right|$; then define $L(s,\pi)\coloneqq(1-\alpha_1q^{-s})^{-1}$ with $\alpha_1$ as above.
	\end{itemize}
	For a character $\xi\colon F^\times\to\CC^\times$, define $L(s,\pi,\xi)\coloneqq L(s,\pi\otimes\xi)$.
\end{definition}
\begin{proposition}
	Let $V$ be a smooth generic irreducible representation of $\op{GL}_2(F_v)$, which we identify with its Kirillov model. Then
	\[\{Z(s,\varphi):\varphi\in V\}=L(s,\pi)\CC[q^s,q^{-s}].\]
	In particular, $Z(s,\varphi)$ admits meromorphic continuation.
\end{proposition}
\begin{proof}[Sketch]
	There is nothing to say in the supercuspidal case. We only sketch the case of $V=\op{Ind}_B^G\chi$ because the others are similar. Because $C_c^\infty(F_v^\times)\subseteq V$ yields $\CC[q^s,q^{-s}]$, it remains to compute $Z\big(s,\left|\cdot\right|^{1/2}\chi_i1_\OO\big)$. This is
	\[\int_\OO\chi_i(y)\left|y\right|^s\,d^\times y=\Bigg(\sum_{k=0}^\infty\chi_i(\varpi)^kq^{-sk}\Bigg)\int_{\OO^\times}\chi_i(t)\,d^\times t.\]
	The integral vanishes if $\chi_i$ fails to be unramified; otherwise, it is a positive constant. The infinite sum is $(1-\chi_i(\varpi)q^{-s})^{-1}$.
\end{proof}
\begin{lemma}[Multiplicity one]
	Let $V$ be a smooth generic irreducible representation of $\op{GL}_2(F_v)$. Embed $F^\times\to T(F)$ by the first coordinate. Given $\chi\colon F^\times\to\CC^\times$, there are at most two $s\in\CC/2\pi i\log q\ZZ$ for which
	\[\dim_{F^\times}(V,\chi\left|\cdot\right|^s)>1.\]
\end{lemma}
\begin{proof}[Sketch]
	Suppose $s\in\CC/2\pi i\log q\ZZ$ features linearly independent functionals $\ell_1,\ell_2\colon V\to\chi\left|\cdot\right|^s$. When restricted to $C_c^\infty(F^\times)$ in the Kirillov model, $\ell_1$ and $\ell_2$ are eigen-distributions and thus unique up to scalar, so there is nonzero $\ell\in\op{span}\{\ell_1,\ell_2\}$ which vanishes on $V(N)$. But then $\ell$ realizes $\chi\left|\cdot\right|^s$ as a one-dimensional quotient of $J(V)$, which occurs for at most two $s$.
\end{proof}
\begin{theorem}[Functional equation]
	Let $V$ be a smooth generic irreducible representation of $\op{GL}_2(F_v)$, which we identify with its Kirillov model. Let $\omega$ be the central character. Then there is a meromorphic function $\gamma(s,\pi,\psi)$ such that
	\[Z\left(1-s,w\varphi,\omega^{-1}\right)=\gamma(s,\pi,\psi)Z(s,\varphi).\]
\end{theorem}
\begin{proof}[Sketch]
	Being meromorphic follows from plugging in a particular $\varphi$ to make $Z(s,\varphi)$ nonzero. Thus, the main issue is independence of $\varphi$: we will check that the two linear functionals
	\[\varphi\mapsto Z(s,\varphi)\qquad\text{and}\qquad\varphi\mapsto Z\left(1-s,w\varphi,\omega^{-1}\right)\]
	are proportional for cofinitely many $s$. Both functionals provide $F^\times$-equivariant maps to $\left|\cdot\right|^{-s+1/2}$, so this follows from the multiplicity one result.
\end{proof}

\subsection{Local Intertwining Operators}
\begin{definition}
	Let $G$ be a split reductive group with Borel $B=TN$. For $w\in W$ and character $\chi\colon T(F_v)\to\CC^\times$, define the intertwining operator $Mw\colon\op{Ind}_B^G\chi\to\op{Ind}_B^G{^w\chi}$ by
	\[M_\chi w(\varphi)(g)\coloneqq\int_{N(F_v)}\varphi(wng)\,dn,\]
	provided the integral converges. We will usually abbreviate $M_\chi w$ to $Mw$ or $M_\chi$ or $M$.
\end{definition}
\begin{remark}
	Direct calculation shows that $M_\chi$ is adjoint to $M_{^{w^{-1}}\overline\chi^{-1}}$. Indeed, for suitable $\varphi_1$ and $\varphi_2$, one has
	\[\int_{K}Mw(\varphi_1)(x)\overline{\varphi_2(x)}\,dx=\int_{K}\int_{N(F_v)}\varphi_1(wnx)\overline{\varphi_2(nx)}\,dn\,dx=\int_{Z(F_v)\backslash G(F_v)}\varphi(wg)\overline{\psi(g)}\,dg.\]
	Now apply $g\mapsto w^{-1}g$ and work backwards.
\end{remark}
\begin{proposition}
	Let $G$ be a split reductive group with Borel $B=TN$. For $w\in W$, the intertwining operator $Mw\colon\op{Ind}_B^G\chi\to\op{Ind}_B^G{^w\chi}$ converges absolutely and uniformly in a suitable cone of $\chi\in X^*(T)_\CC$.
\end{proposition}
\begin{proof}[Sketch]
	We only argue this for $G=\op{GL}_2$ and $v$ finite. Then we are interested in the convergence of
	\[Mw(\varphi)(g)=\int_{F_v}\varphi\left(w\begin{bmatrix}
		1 & x \\ & 1
	\end{bmatrix}g\right)\,dx.\]
	The region $\left|x\right|<1$ is compact and gives no issue. The region $\left|x\right|\ge1$ is bounded using the ``by-hand'' Iwasawa decomposition
	\[\begin{bmatrix}
		& -1 \\ 1
	\end{bmatrix}\begin{bmatrix}
		1 & x \\ & 1
	\end{bmatrix}=\begin{bmatrix}
		1/x & -1 \\ & x
	\end{bmatrix}\begin{bmatrix}
		1 \\ 1/x & 1
	\end{bmatrix}.\]
	One then factors out $\int_{\left|x\right|\le1}\left|x\right|^{2\sigma-2}\,dx$, where $\sigma\in\RR$ is chosen so that $\left|\chi_1\chi_2^{-1}\right|=\left|\cdot\right|^\sigma$. Thus, we receive convergence when $\sigma>1/2$.
\end{proof}
\begin{definition}
	Define the vector bundle $\int_{X^*(T)_\CC}^\oplus\op{Ind}_B^G\chi\,d\chi$ over $X^*(T)_\CC$ to be the locally trivial fiber bundle whose fiber over $\chi$ is $\op{Ind}_B^G\chi$. Here the trivialization $\op{Ind}_B^G\chi\to\op{Ind}_B^G\chi\left|\cdot\right|^s$ sends $\varphi$ to
	\[\varphi_s(g)\coloneqq m_B(g)^{s\delta}\varphi(g),\]
	where $m_B(g)$ is the Levi part of the Iwasawa decomposition of $g$.
\end{definition}
\begin{definition}
	Take $G=\op{GL}_2$. Fix a character $\chi\colon T(F_v)\to\CC^\times$. For a Schwartz function $\Phi\colon F_v\times F_v\to\CC$, define
	\[\varphi_\Phi(g;\chi)\coloneqq\frac{\chi_1\left|\cdot\right|^{1/2}(\det g)}{L\left(1,\chi_1\chi_2^{-1}\right)}\int_{F_v^\times}\Phi(te_2g)\chi_1\chi_2^{-1}(t)\left|t\right|\,d^\times t.\]
	The integral is defined via analytic continuation by realizing that it is a local $Z$-function in the sense of Tate's thesis. Note $\varphi_\Phi$ defines a smooth global section of $\int_{X^*(T)_\CC}\op{Ind}_B^G\chi\,d\chi$.
\end{definition}
\begin{lemma}
	Take $G=\op{GL}_2$. Every $\varphi\in\op{Ind}_B^G\chi$ admits some Schwartz $\Phi_\varphi\colon F_v\times F_v\to\CC$ for which $\varphi=\varphi_{\Phi_\varphi}$.
\end{lemma}
\begin{proof}[Sketch]
	Define
	\[\overline\Phi_\varphi(v)\coloneqq\begin{cases}
		\chi_1^{-1}(\det g)f(g) & \text{if }v=e_2g, \\
		0 & \text{otherwise}.
	\end{cases}\]
	This is well-defined. A direct calculation shows that $\varphi\simeq\varphi_{\overline\Phi_\varphi}$, where the constant is given by the measure of $\OO^\times$ and $L\left(1,\chi_1\chi_2^{-1}\right)$. Rescaling provides $\Phi_\varphi$.
\end{proof}
\begin{proposition}
	Take $G=\op{GL}_2$, and fix a character $\chi\colon T(F_v)\to\CC^\times$ in a suitable cone. For Schwartz $\Phi\colon F_v\times F_v\to\CC$, one has
	\[Mw(\varphi_\Phi)=\left(\frac{L(0,\chi_1\chi_2^{-1})}{\varepsilon(\chi_1\chi_2^{-1};\psi)L(1,\chi_1\chi_2^{-1})}\right)\cdot\chi_1\chi_2(-1)\cdot\frac{\chi_2\left|\cdot\right|^{1/2}(\det g)}{L\left(1,\chi_2\chi_1^{-1}\right)}\int_{F_v^\times}\mc F_\psi\Phi(te_2g)\chi_1^{-1}\chi_2(t)\left|t\right|\,d^\times t,\]
	where $\mc F_\psi\Phi$ is the Fourier transform with respect to $\psi\colon F_v\to\CC^\times$.
\end{proposition}
\begin{proof}[Sketch]
	This is a long calculation. We remark that
	\[\mc F_\psi\Phi(w)=\int_{F_v^2}\Phi(v)\psi\left(\det\begin{bmatrix}
		v & w
	\end{bmatrix}\right)\,dv.\]
	Now, the main point is to use Tate's thesis to ``flip'' the character in the integral from $\chi_1\chi_2^{-1}$ to $\chi_2^{-1}\chi_1$. In particular, one must compute the Fourier transform of $t\mapsto\Phi(te_2g)$, which is $t\mapsto\mc F_\psi\Phi(te_2g)$.
\end{proof}
\begin{theorem}
	Let $G$ be a split reductive group with Borel $B=TN$. For $w\in W$, the intertwining operator $Mw$ admits meromorphic continuation to a meromorphic operator
	\[\int_{X^*(T)_\CC}^\oplus\op{Ind}_B^G\chi\,d\chi\to\int_{X^*(T)_\CC}^\oplus\op{Ind}_B^G{^w\chi}\,d\chi.\]
\end{theorem}
\begin{proof}[Sketch]
	We only argue this for $G=\op{GL}_2$. Simply define
	\[Mw(\varphi)\coloneqq\left(\frac{L(0,\chi_1\chi_2^{-1})}{\varepsilon(\chi_1\chi_2^{-1};\psi)L(1,\chi_1\chi_2^{-1})}\right)\cdot\chi_1\chi_2(-1)\cdot\frac{\chi_2\left|\cdot\right|^{1/2}(\det g)}{L\left(1,\chi_2\chi_1^{-1}\right)}\int_{F_v^\times}\mc F_\psi\Phi_\varphi(te_2g)\chi_1^{-1}\chi_2(t)\left|t\right|\,d^\times t,\]
	where as usual the right-hand side is defined by meromorphic continuation from Tate's thesis. To check that $Mw$ is a meromorphic operator, it is enough to check that, for any flat sections $f_s$ and $g_s$, the function $s\mapsto\langle Mw(f_s),g_s\rangle$ is meromorphic. This is direct from the definition.
\end{proof}
\begin{notation}
	For later use, let $Rw$ be equal to the operator
	\[Rw(\varphi)\coloneqq\chi_1\chi_2(-1)\cdot\frac{\chi_2\left|\cdot\right|^{1/2}(\det g)}{L\left(1,\chi_2\chi_1^{-1}\right)}\int_{F_v^\times}\mc F_\psi\Phi_\varphi(te_2g)\chi_1^{-1}\chi_2(t)\left|t\right|\,d^\times t.\]
	Roughly speaking, $Rw$ is a ``normalized'' variant of $Mw$, and it admits holomorphic continuation.
\end{notation}
\begin{example}
	If $\chi$ is unramified, let $\varphi_\chi\in\op{Ind}_B^G\chi$ denote the unramified vector satisfying $\varphi_\chi(1)=1$. Then $Rw(\varphi_\chi)\in\op{Ind}_B^G{^w\chi}$ is also $K$-invariant, so it is a scalar multiple of the corresponding unramified vector. A local calculation (with $Mw$) given by plugging in $1$ shows that $Rw(\varphi_\chi)=\varphi_{^w\chi}$.
\end{example}
\begin{remark}
	Because the Fourier transform is self-dual, one finds that $Rw\circ Rw=1$.
\end{remark}
\begin{remark}
	If $\eta$ is unitary, then $R_\eta w$ is unitary, which we see by using $Rw\circ Rw=1$ and the adjoint property for $Mw$.
\end{remark}
\begin{remark}
	Consider the case where $\chi_1\chi_2^{-1}=\left|\cdot\right|$ so that $\op{Ind}_B^G{^w\chi}$ admits a one-dimensional subspace. Set $\rho\coloneqq\chi_1\left|\cdot\right|^{-1/2}$; then the one-dimensional subspace is $g\mapsto\chi_v(\det g)$. Now, $Rw\colon\op{Ind}_B^G\chi\to\op{Ind}_B^G{^w\chi}$ is nonzero, so its image must be a quotient of $\op{Ind}_B^G\chi$ and a subrepresentation of $\op{Ind}_B^G{^w\chi}$. Passing to the Jacquet module, we see that $Rw$ must be a scalar times
	\[\op{Ind}_B^G\chi\onto{\rho\circ\det}\subseteq\op{Ind}_B^G{^w\chi}.\]
	In particular, $Rw(\varphi)=c(\varphi)({\rho\circ\det})$ for some scalar $c(\varphi)$, and invariance properties of $c$ force the map $c\colon\op{Ind}_B^G\chi\to\CC$ to be $\left\langle-,\rho\circ\det\right\rangle$. It follows that the matrix coefficient of $Rw$ is
	\[\langle Rw(\varphi_1),\varphi_2\rangle\simeq\langle\varphi_1,{\rho\circ\det}\rangle\overline{\langle\varphi_2,{\rho\circ\det}\rangle},\]
	where the implied constant could be computed explicitly.
\end{remark}

\section{The Cuspidal Spectrum}

\subsection{The Cuspidal Subspace is Closed}
\begin{definition}
	Some $\varphi\in L^2([G])$ is \textit{cuspidal} if and only if
	\[\int_{[N]}\varphi(ng)\,dn=0\]
	for almost every $g\in A_G\backslash G(\AA_F)$ and each unipotent radical $N$ of a proper parabolic subgroup $P$ of $G$. We define $L^2_{\mathrm{cusp}}([G])$ to be the subspace of cuspidal functions.
\end{definition}
\begin{example}
	If $G$ admits no proper parabolic subgroups, then $L^2([G])=L^2_{\mathrm{cusp}}([G])$.
\end{example}
\begin{definition}[period integral]
	Fix a subgroup $H\subseteq G$ and a character $\chi\colon A_GH(F)\backslash A_GH(\AA_F)\to\CC$. For continuous $\varphi\colon[G]\to\CC$, we define the \textit{period}
	\[\mc P_\chi(\varphi)\coloneqq\int_{A_GH(F)\backslash A_GH(\AA_F)}\varphi(h)\chi^{-1}(h)\,dh\]
	as long as the integral absolutely converges.
\end{definition}
\begin{notation}
	Fix a unipotent subgroup $H\subseteq G$ and a character $\chi\colon A_GH(F)\backslash A_GH(\AA_F)\to\CC$. We define the temporary notation
	\[L^2([G])(\chi)\coloneqq\left\{\varphi\in L^2([G]):\mc P_\chi(Rg(\varphi))=0\text{ for almost all }g\right\}.\]
\end{notation}
\begin{example}
	Note that $L^2_{\mathrm{cusp}}([G])$ is
	\[\bigcap_{\substack{\text{parabolic }P\\P\ne G\\P=MN}}L^2([G])(1_N).\]
\end{example}
\begin{remark}
	The condition $\mc P_\chi(Rg(\varphi))=0$ for almost every $g\in A_G\backslash G(\AA_F)$ is equivalent to the condition $\mc P_\chi(Rf(\varphi))=0$ for all $f\in C_c^\infty(A_G\backslash G(\AA_F))$. Indeed, by rearranging the order of integration,
	\[\int_{A_G\backslash G(\AA_F)}f(g)\mc P_\chi(Rg(\varphi))\,dg=\mc P_\chi(Rf(\varphi)).\]
	Thus, $\varphi\in L^2([G])(\chi)$ does imply $\mc P_\chi(Rf(\varphi))=0$. For the converse, do an approximation argument.
\end{remark}
\begin{definition}[automorphic kernel]
	For $f\in C_c^\infty(A_G\backslash G(\AA_F))$, define $K_f\in C^\infty([G]\times[G])$ by
	\[K_f(x,y)\coloneqq\sum_{\gamma\in G(F)}f\left(x^{-1}\gamma y\right).\]
	The sum converges because it is finite.
\end{definition}
\begin{remark}
	By unfolding, the integral operator given by the kernel $K_f$ on $L^2([G])$ is exactly $Rf$.
\end{remark}
\begin{proposition}
	The cuspidal subspace $L^2_{\mathrm{cusp}}([G])\subseteq L^2([G])$ is closed.
\end{proposition}
\begin{proof}[Sketch]
	It is enough to show the operators $\varphi\mapsto\mc P_\chi(Rf(\varphi))$ are continuous, where $f\in C_c^\infty(A_G\backslash G(\AA_F))$. By unfolding, we see that
	\[\mc P_\chi(Rf(\varphi))=\int_{[G]}\operatorname{Pe}_{f,\chi\delta_G}(g)\varphi(g)\,dg,\]
	where
	\[\operatorname{Pe}_{f,\chi}(g)\coloneqq\int_{A_GH(F)\backslash A_GH(\AA_F)}K_f(h,g)\chi^{-1}(h)\,dh.\]
	It is thus enough to show that $\mathrm{Pe}_{f,\chi}\in L^2([G])$, which is a direct calculation.
\end{proof}

\subsection{The Cuspidal Subspace is Discrete}
\begin{notation}
	Suppose $G$ is reductive. For $\alpha\in\Delta$, let $\Phi_\alpha^+$ denote the collection of positive roots which have $\alpha$ in their support. Let $N(\Phi_\alpha^+)$ denote the product of the corresponding root spaces.
\end{notation}
% \begin{lemma}
% 	Suppose $F$ is a function field and that $G$ is reductive. 
% \end{lemma}
% \begin{proof}[Sketch]
% 	This is a rather intricate induction on root spaces. By using the upper central series, one can choose a filtration $\{N_i\}_i$ of $N$ so that $N_0=N$ and $N_1=N(\Phi_\alpha^+)$ and $N_m=1$ for $m$ large enough, and each quotient $N_i/N_{i+1}$ is isomorphic to a root space $\mathbb G_a$. (In fact, $[N,N_i]\subseteq N_{i+1}$.) The moral is that the filtration allows us to pass to the case where $N(\Phi_\alpha^+)$ is isomorphic to $\mathbb G_a$, and then we conclude by using the compactness of $F\backslash\AA_F$.
% \end{proof}
\begin{theorem}[Harder]
	Suppose $F$ is a function field and that $G$ is reductive. Then any smooth $\varphi\in L^2_{\mathrm{cusp}}([G])$ is compactly supported.
\end{theorem}
\begin{proof}[Sketch]
	As is typical in proofs involving Siegel sets, we have a combinatorial step, and then we apply it to the Siegel sets of interest.
	\begin{enumerate}
		\item Choose a compact open subgroup $K'\subseteq G(\AA_F)$. Then for $t\gtrsim_{K'}0$ and $h\in\mf S(t)$, we claim that
		\[N(\Phi_\alpha^+)(F)\cdot\left(N(\Phi_\alpha^+)\cap hK'h^{-1}\right)\stackrel?=N(\Phi_\alpha^+)(\AA_F).\]
		This is a rather intricate induction on root spaces. By using the upper central series, one can choose a filtration $\{N_i\}_i$ of $N$ so that $N_0=N$ and $N_1=N(\Phi_\alpha^+)$ and $N_m=1$ for $m$ large enough, and each quotient $N_i/N_{i+1}$ is isomorphic to a root space $\mathbb G_a$. (In fact, $[N,N_i]\subseteq N_{i+1}$.) The moral is that the filtration allows us to pass to the case where $N(\Phi_\alpha^+)$ is isomorphic to $\mathbb G_a$, and then we conclude by using the compactness of $F\backslash\AA_F$.
		\item We now conduct the proof. By formalism involving Siegel sets, it is enough to show that $\varphi(g)=0$ for $g\in\mf S(t)$, where $t\gtrsim_\varphi0$. By smoothness, $\varphi$ is invariant under some compact open $K'\subseteq G(\AA_F)$. Then choose $t$ as in the previous step. Now, for any $h\in\mf S(t)$ and $g\in A_GG(\AA_F)^1\backslash G(\AA_F)$, we see that the function $N(\Phi_\alpha^+)(\AA_F)\to\CC$ defined by
		\[x\mapsto\varphi(xhg_i)\]
		is invariant under $N(\Phi_\alpha^+)(\AA_F)$ (by the lemma), so they must vanish uniformly because $\varphi$ is cuspidal.
		\qedhere
	\end{enumerate}
\end{proof}
\begin{corollary}
	Suppose $F$ is a function field and that $G$ is reductive. For $f\in C_c^\infty(A_G\backslash G(\AA_F))$, the operator $Rf$ on $L^2_{\mathrm{cusp}}([G])$ is trace class.
\end{corollary}
\begin{proof}[Sketch]
	By the Dixmier--Malliavin theorem, it is enough to show that $Rf$ is Hilbert--Schmidt.
	
	Now, the idea is that the continuity of $K_f$ means that $Rf$ is Hilbert--Schmidt when restricted to $L^2(\Omega)$ for any compact $\Omega\subseteq[G]$. To be precise, choose compact open $K\subseteq G(\AA_F)$ for which $f\in C_c^\infty(G(\AA_F))^K$. Then $Rf$ projects $L^2_{\mathrm{cusp}}([G])$ onto the $K$-invariants, so it is enough to check that $Rf$ is Hilbert--Schmidt on a basis of $L^2_{\mathrm{cusp}}([G])^K$. By Harder's theorem, all such functions are compactly supported in some $\Omega\subseteq G(\AA_F)$, so we are done because $Rf|_{L^2(\Omega)}$ is Hilbert--Schmidt.
\end{proof}
\begin{theorem}
	Suppose $F$ is a number field and that $G$ is reductive. Choose $f\in C_c^\infty(A_G\backslash G(\AA_F))$ and $B>0$ and $t>0$. For any $\varphi\in L^2_{\mathrm{cusp}}([G])$ and $x\in\mf S(t)$, we have
	\[\left|Rf(\varphi)(x)\right|\lesssim\rho(s_x)^{-B}\norm\varphi_2.\]
\end{theorem}
\begin{proof}[Sketch]
	This is a hard theorem! By multiplying the inequality by itself, we may replace $\rho$ by $\alpha$, where $\alpha\in\Delta$. By restricting scalars, we may as well take $F=\QQ$, but we will avoid using this hypothesis.
	\begin{enumerate}
		\item We separate frequencies. By unfolding,
		\[Rf(\varphi)(x)=\int_{A_GN(F)\backslash G(\AA_F)}\Bigg(\sum_{\gamma\in N(F)}f\left(x^{-1}\gamma y\right)\Bigg)\varphi(y)\,dy.\]
		Fix a nontrivial character $\psi\colon[\mathbb G_a]\to S^1$. By Poisson summation,
		\[Rf(\varphi)(x)=\int_{A_GN(F)\backslash G(\AA_F)}\Bigg(\sum_{\eta\in\mf n(\QQ)}\underbrace{\int_{\mf n(\AA_F)}f\left(x^{-1}\exp(n)y\right)\psi(\langle\eta,n\rangle)\,dn}_{\widehat f_\eta\left(x^{-1}\exp(n)y\right)\coloneqq}\Bigg)\varphi(y)\,dy.\]
		
		\item We now apply the cuspidal hypothesis to get rid of some of the frequencies: we claim that we can remove the $\eta\in\mf n(\QQ)$ for which $\langle\eta,\mf n(\Phi_\alpha^+)\rangle=0$. Well, for such $\eta$, an application of the Baker--Campbell--Hausdorff formula shows that $y\mapsto\widehat f_\eta\left(x^{-1}\exp(n)y\right)$ is left-invariant by $N(\Phi_\alpha^+)(\AA_F)$. Such a variable change is then able to produce an integral over $[N(\Phi_\alpha^+)]$ inside the integral for $Rf(\varphi)(x)$ at these frequencies.

		\item There is some very delicate combinatorics of Siegel sets that we will mostly avoid, but let's give the output. We may assume that $f$ factors as $f_\infty\otimes f^\infty\in C_c^\infty(A_G\backslash G(F_\infty))\otimes C_c^\infty(G(\AA_F^\infty))$. By keeping track of supports (and then the Fourier transforms), we find that
		\[\int_{\mf n(\AA_F^\infty)}f^\infty\left(x^{\infty,-1}\exp(n)y^\infty\right)\psi(\langle\eta,n\rangle)\,dn\lesssim1_{\Omega_{\mf n}^\infty}(\eta)\]
		for some compact $\Omega_{\mf n}^\infty$. More combinatorics of the supports shows that
		\[\int_{\mf n(F_\infty)}f_\infty\left(x_\infty^{-1}\exp(n)y_\infty\right)\psi(\langle\eta,n\rangle)\,dn\lesssim\delta_P(s_x)\norm{\op{Ad}^\lor\left(s_x^{-1}\right)\eta}^{-B}.\]
		
		\item It remains to lower-bound $\norm{\op{Ad}^\lor\left(s_x^{-1}\right)\eta}$ when $\eta\in\mf(\AA_F)\cap\Omega_{\mf n}^\infty$; in other words, $\eta$ lives in a lattice. By considering the diagonal action of the adjoint action on the root spaces, we find $\norm{\op{Ad}^\lor\left(s_x^{-1}\right)\eta}\gtrsim_t\norm\eta_{\mf n}$; alternatively, computing this action on a single root $\lambda\in\Phi_\alpha^+$ gives $\norm{\op{Ad}^\lor\left(s_x^{-1}\right)\eta}\gtrsim\lambda(s_x)\norm{\eta_\lambda^\lor}$, and this norm is constantly lower-bounded by being in a lattice. Combining the two bounds gives
		\[\int_{\mf n(F_\infty)}f_\infty\left(x_\infty^{-1}\exp(n)y_\infty\right)\psi(\langle\eta,n\rangle)\,dn\lesssim\delta_P(s_x)\alpha(s_x)^{-B}\norm\eta^{-B}1_{\Omega_{\mf n}^\infty}.\]
		Summing and integrating bounds $\left|Rf(\varphi)(x)\right|$ in terms of an $L^1$ norm of $\varphi$, and we conclude by H\"older's inequality.
		\qedhere
	\end{enumerate}
\end{proof}
\begin{corollary}
	Suppose $F$ is a number field and that $G$ is reductive. For $f\in C_c^\infty(A_G\backslash G(\AA_F))$, the operator $Rf$ on $L^2_{\mathrm{cusp}}([G])$ is trace class.
\end{corollary}
\begin{proof}[Sketch]
	By the Dixmier--Malliavin theorem, it is enough to show that $Rf$ is Hilbert--Schmidt. The theorem implies that $\norm{Rf(\varphi)}_\infty\lesssim_f\norm\varphi_2$. The idea is to realize $Rf$ as the integral operator of a kernel $K\in L^2([G]\times[G])$. For each $x\in[G]$, we may use the Riesz representation theorem to define $K_x\in L^2([G])$ so that
	\[Rf(\varphi)(x)=\int_{[G]}K_x(y)\varphi(y)\,dy.\]
	Now, define $K\colon[G]\times[G]\to\CC$ (almost everywhere) by $K(x,y)\coloneqq K_x(y)$. By construction $Rf$ is given by the kernel $K$, so it remains to check $\norm K_2$. Well, the bound from the theorem gives $\norm{K_x}_2\lesssim_f1$, so one finds that $\norm K_2$ is finite.
\end{proof}

\section{Decomposition of the Spectrum}

\subsection{Eisenstein Series}
\begin{definition}[discrete spectrum]
	Let $G$ be a reductive group. Let $L^2_{\mathrm{disc}}([G])$ be the largest closed subspace which decomposes into a Hilbert direct sum of irreducible subrepresentations.
\end{definition}
\begin{remark}
	Note that $L^2_{\mathrm{disc}}([G])$ exists by Zorn's lemma. It contains $L^2_{\mathrm{cusp}}([G])$.
\end{remark}
\begin{notation}
	Let $G$ be a reductive group with parabolic subgroup $P=MN$. Choose a discrete subrepresentation $\sigma\subseteq L^2_{\mathrm{disc}}([M])$. Following Arthur, we define the vector space $\op{Ind}_P^G\sigma$ to consist of the measurable functions $N(\AA_F)A_MM(\AA_F)\backslash G(\AA_F)\to\CC$ for which
	\begin{itemize}
		\item the map $m\mapsto\varphi(mx)$ is in $\sigma$ for all $x$, and
		\item the norm
		\[\norm\varphi^2\coloneqq\int_K\int_{A_MM(F)\backslash M(\AA_F)}\left|\varphi(mk)\right|^2\,dm\,dk\]
		is finite.
	\end{itemize}
	Given $\lambda\in\mf a^*_{M\CC}$, we define the representation $I(\sigma,\lambda)$ on the space $\op{Ind}_P^G\sigma$ as
	\[R_\lambda g(\varphi)(x)\coloneqq\varphi(xg)\exp(\langle H_P(xg)-H_P(x),\lambda+\rho_P\rangle).\]
\end{notation}
\begin{remark}
	One could alternatively view $I(\sigma,\lambda)$ as $\op{Ind}_P^G\left(\sigma\otimes e^{\langle H_P(-),\lambda\rangle}\right)$.
\end{remark}
\begin{remark}
	A direct calculation shows that $I(\sigma,\lambda)$ is unitary if $\lambda\in i\mf a_M^*$.
\end{remark}
\begin{notation}
	Let $G$ be a reductive group with parabolic subgroup $P=MN$. Choose a discrete subrepresentation $\sigma\subseteq L^2_{\mathrm{disc}}([M])$. Define $\op{Ind}_P^G(\sigma)^\circ$ to consist of the smooth $K$-finite vectors of $\op{Ind}_P^G(\sigma)$.
\end{notation}
\begin{definition}[Eisenstein series]
	Let $G$ be a reductive group with parabolic subgroup $P=MN$. Choose an irreducible $\sigma\subseteq L^2_{\mathrm{disc}}([M])$. For $\varphi\in\sigma$, we ddefine the \textit{Eisenstein series}
	\[E(x,\varphi,\lambda)\coloneqq\sum_{\gamma\in P(F)\backslash G(F)}\varphi(\gamma x)e^{\langle H_P(\gamma x),\lambda+\rho_P\rangle}.\]
\end{definition}
\begin{proposition}
	Let $G$ be a reductive group with parabolic subgroup $P=MN$. Choose an irreducible $\sigma\subseteq L^2_{\mathrm{cusp}}([M])$. Then $E(x,\varphi,\lambda)$ converges absolutely and uniformly on compacts of $G(\AA_F)$ if $\Re\lambda$ lives in a suitable Weyl cone.
\end{proposition}
\begin{proof}[Sketch]
	Because $\sigma$ is cuspidal, $\varphi$ is rapidly decreasing and in particular bounded. Thus, it remains to bound
	\[\sum_{\gamma\in P(F)\backslash G(F)}\left|m_P(\gamma x)^{\lambda+\rho_P}\right|.\]
	Let $U\subseteq G(\AA_F)^1$ be a small open relatively compact neighborhood of $1$. Then uniform continuity allows us to replace each summand with an integral over $U$, so we want to bound
	\[\int_{P(F)\backslash G(F)V}\left|m_P(yx)^{\lambda+\rho_P}\right|\,dy\]
	after unfolding. Because $P\backslash G$ is projective, $P(\AA_F)\backslash G(\AA_F)$ is compact, so we need to bound the integral over $P(F)\backslash P(\AA_F)$. Unipotent radicals are compact, so we may even bound the integral over $M(F)\backslash M(\AA_F)$. By passing to a Siegel set, we may go to the integral over $A_M(F)\backslash A_M(\AA_F)$. Thus, we are basically integrating over a split torus, which will be bounded for suitable exponends $\lambda$.
\end{proof}
\begin{remark}
	A more careful analysis of how $x$ varies over a Siegel set can show that $E(x,\varphi,\lambda)$ has moderate growth for these $\lambda$.
\end{remark}
\begin{theorem}
	Let $G$ be a reductive group with parabolic subgroup $P=MN$. Choose an irreducible $\sigma\subseteq L^2_{\mathrm{cusp}}([M])$. For $\varphi\in\sigma$, the Eisenstein series $E(x,\varphi,\lambda)$ admits meromorphic continuation to all $\lambda$.
\end{theorem}
\begin{proof}[Sketch]
	We only argue this for $G=\op{GL}_2$, and we will do it later.
\end{proof}

\subsection{Intertwining Operators}
\begin{notation}
	Let $M'\subseteq M$ be Levi subgroups of $G$. Then define $\mf a_M^*\coloneqq X^*(M)_\RR$ and $\mf a_M$ as its dual, and define $\mf a_{M'}$ similarly. Then $\mf a_{M'}^{M}$ is the kernel of the projection $\mf a_{M'}\to\mf a_M$ (induced by the dual of restriction). 
\end{notation}
\begin{definition}
	Let $G$ be a reductive group, and let $M$ and $M'$ be standard Levi subgroups of $G$.
	\begin{itemize}
		\item Define $W^{M,M'}$ to consist of the $w\in W$ for which $w\Phi^+(M)\subseteq\Phi^+(G)$ and $w^{-1}\Phi^+(M')\subseteq\Phi^+(G)$.
		\item Define $W(M,M')$ to consist of the $w\in W^{M,M'}$ for which $wMw^{-1}\subseteq M'$.
	\end{itemize}
\end{definition}
\begin{remark}
	It turns out that $W^{M_I,M_J}$ uniquely represents the coset space $W_J\backslash W/W_I$ by elements of minimal length.
\end{remark}
\begin{remark}
	The Bruhat decomposition upgrades to
	\[G=\bigsqcup_{w\in W_I\backslash W/W_J}P_IwP_J.\]
\end{remark}
\begin{remark}
	An explicit calculation on the Lie algebra (with root spaces) shows that $w^{-1}P_Jw\cap M_I$ is a standard parabolic subgroup of $M_I$ with Levi decomposition
	\[w^{-1}P_Jw\cap M_I=\left(w^{-1}M_Jw\cap M_I\right)\left(w^{-1}N_Jw\cap M_I\right).\]
\end{remark}
\begin{definition}[associate]
	Two parabolic subgroups $P$ and $P'$ of $G$ are \textit{associate} if and only if their Levi subgroups are conjugate by $G$.
\end{definition}
\begin{notation}
	Let $M$ and $M'$ be standard Levi subgroups of a reductive group $G$. Let $W(\mf a_M,\mf a_{M'})$ consist of those $w\in W$ inducing an isomorphism $\mf a_M\to\mf a_{M'}$.
\end{notation}
\begin{remark}
	Let $P_0=M_0N_0$ be a minimal parabolic. If $P$ and $P'$ are stadnard associate with Levi subgroups $M$ and $M'$, one can show that there is $w\in W$ such that the automorphism $w\colon\mf a_{M_0}\to\mf a_{M_0}$ which restricts to an isomorphism $\mf a_M\to\mf a_{M'}$. (Such a $w$ also conjugates the Levi subgroups to each other.)
\end{remark}
\begin{definition}
	Let $P$ and $P'$ be associate standard parabolic subgroups of a reductive group $G$ with $P=MN$ and $P'=M'N'$, and choose $w\in W(\mf a_M,\mf a_{M'})$. For irreducible $\sigma\subseteq L^2_{\mathrm{disc}}([M])$ and $\lambda\in\mf a_{M,\CC}$, define $M_\lambda w\colon I(\sigma,\lambda)\to I(\sigma,w\lambda)$ by having $M_\lambda w(\varphi)(x)$ equal
	\[\int_{\left(N'(\AA_F)\cap wN(\AA_F)w^{-1}\right)\backslash N'(\AA_F)}\varphi\left(w^{-1}nx\right)\exp(\langle H_P(w^{-1}nx),\lambda+\rho_P\rangle+\langle H_{P'}(x),-w\lambda+\rho_{P'}\rangle)\,dn,\]
	provided that the integral converges.
\end{definition}
\begin{remark}
	By folding the integral into a sum over $F$-points, one can turn this into an integral over an Eisenstein series over a compact group. Thus, for $\Re\lambda$ in a suitable cone, the integral converges absolutely.
\end{remark}
\begin{notation}
	Fix a reductive group $G$. Let $\mc P$ be an association class of standard parabolic subgroups of $G$. Then let $\mc H_{\mc P}$ denote the Hilbert space of tuples of measurable functions $\{F_P\}_{P\in\mc P}$ such that
	\begin{itemize}
		\item $F_P\colon i\mf a_{M\RR}^{G*}\to\op{Ind}_P^GL^2_{\mathrm{disc}}([M])$ satisfies $F_{P'}(w\lambda)=M_\lambda wF_P(\lambda)$ if $w\in W(\mf a_m,\mf a_{M'})$, and
		\item The norm
		\[\norm F^2\coloneqq\sum_{P\in\mc P}n_{P}^{-1}\int_{i\mf a^{G*}_{M\RR}}\norm{F_P(\lambda)}^2\,d\lambda\]
		is finite, where $n_P\coloneqq\sum_{P'\in\mc P}\#W(\mf a_M,\mf a_{M'})$.
	\end{itemize}
\end{notation}
\begin{notation}
	Fix a reductive group $G$, and let $\mc P$ be an association class of standard parabolic subgroups. Define $E_{\mc P}\colon\mc H_P\to L^2([G])$ by
	\[E_{\mc P}(F)(x)\coloneqq\sum_{P\in\mc P}n_P^{-1}\int_{i\mf a_{M\RR}^{G*}}E(x,F_P(\lambda),\lambda)\,d\lambda.\]
	The image of $E_{\mc P}$ is denoted $L_{\mc P}([G])$.
\end{notation}
\begin{theorem}[Langlands]
	Fix a reductive group $G$. Then
	\[L^2([G])=\bigoplus_{\mc P}L^2_{\mc P}([G]),\]
	where the sum is over all association classes (of standard parabolics).
\end{theorem}
\begin{proof}[Sketch]
	This theorem is extremely hard. We will say nothing about it.
\end{proof}

% \subsection{Intertwining Operators for \texorpdfstring{$\op{GL}_2$}{GL(2)}}


\subsection{Preliminaries for \texorpdfstring{$\op{GL}_2$}{GL(2)}}
\begin{notation}
	Throughout, we fix a unitary central character $\omega\colon Z(\AA_F)\to\CC^\times$, which we assume has $\omega|_{F_\infty^+}=1$.
\end{notation}
\begin{definition}[$P$-series]
	For $f\colon N(\AA_F)P(F)(Z(\AA_F),\omega)\backslash G(\AA_F)\to\CC$, define the \textit{$P$-series}
	\[F_f(x)\coloneqq\sum_{\gamma\in P(F)\backslash G(F)}f(\gamma x).\]
\end{definition}
\begin{remark}
	If $f$ is compactly supported modulo $N(\AA_F)Z(\AA_F)P(F)$, then the series is finite.
\end{remark}
\begin{lemma}
	Suppose $\psi\in C_c^\infty(G(F)\backslash G(\AA_F),\omega)$, then
	\[\langle\psi,F_f\rangle=\int_{P(F)Z(\AA_F)N(\AA_F)\backslash G(\AA_F)}\psi_N(x)\overline{f(x)}\,dx.\]
\end{lemma}
\begin{proof}[Sketch]
	This is a direct calculation. The idea is to fold $\langle\psi,F_f\rangle$ into
	\[\int_{\overline G(F)\backslash\overline G(\AA_F)}\Bigg(\sum_{\gamma\in P(F)\backslash G(F)}\psi(\gamma g)\overline{f(\gamma g)}\Bigg)\,dg=\int_{P(F)Z(\AA_F)\backslash G(\AA_F)}\psi(x)\overline f(x)\,dx,\]
	from which the lemma follows.
\end{proof}
\begin{remark}
	It follows that the $P$-series are dense in $L^2_{\mathrm{cusp}}(G(F)\backslash G(\AA_F),\omega)^\perp$.
\end{remark}
\begin{notation}
	For $f\colon N(\AA_F)P(F)(Z(\AA_F),\omega)\backslash G(\AA_F)\to\CC$, define
	\[\widetilde f(g)\coloneqq\int_{N(\AA_F)}f(wng)\,dn.\]
	We may also denote this integral operator by $Mw$.
\end{notation}
\begin{remark}
	The Bruhat decomposition (and unfolding) implies $F_{f,N}(g)=f(g)+\widetilde f(g)$, so
	\[\langle F_{f_1},F_{f_2}\rangle=\int_{P(F)Z(\AA_F)N(\AA_F)\backslash G(\AA_F)}f_1(x)\overline{f_2(x)}\,dx+\int_{P(F)Z(\AA_F)N(\AA_F)\backslash G(\AA_F)}\widetilde f_1(x)\overline{f_2(x)}\,dx.\]
\end{remark}
\begin{notation}
	For $f\colon N(\AA_F)P(F)(Z(\AA_F),\omega)\backslash G(\AA_F)\to\CC$, we define the \textit{Mellin transform}
	\[\widehat f(g;s)\coloneqq\int_{F_\infty^+}f\left(\begin{bmatrix}
		t \\ & 1
	\end{bmatrix}g\right)\left|t\right|^{-s-1/2}\,d^\times t.\]
\end{notation}
\begin{remark}
	Note $\widehat f(-;\lambda)$ descends to $N(\AA_F)T(F)(Z(\AA_F),\omega)\backslash G(\AA_F)$. Thus, it makes sense to define the inner product
	\[\langle\widehat f_1(s),\widehat f_2(-\overline s)\rangle\coloneqq\int_{K}\int_{F^\times\backslash\AA_F^1}\widehat f_1(\op{diag}(a,1)k;s)\overline{\widehat f_2(\op{diag}(a,1)k;-\overline s)}\,da\,dk.\]
\end{remark}
\begin{proposition}
	For $f_1,f_2\colon N(\AA_F)P(F)(Z(\AA_F),\omega)\backslash G(\AA_F)\to\CC$ and $x>1/2$, one has
	\[\langle F_{f_1},F_{f_2}\rangle=\frac1{2\pi i}\int_{\op{Re}\lambda=x}\left\langle\widehat f_1(\lambda),\widehat f_2(-\overline\lambda)\right\rangle\,d\lambda+\frac1{2\pi i}\int_{\op{Re}\lambda=x}\left\langle M_\lambda w(\widehat f_1)(\lambda),\widehat f_2(\overline\lambda)\right\rangle\,d\lambda.\]
\end{proposition}
\begin{proof}[Sketch]
	Applying the Iwasawa decomposition and then Mellin inversion yields
	\[\int_{P(F)Z(\AA_F)N(\AA_F)\backslash G(\AA_F)}f_1(x)\overline{f_2(x)}\,dx=\frac1{2\pi i}\int_{\op{Re}\lambda=x}\left\langle\widehat f_1(\lambda),\widehat f_2(-\overline\lambda)\right\rangle\,d\lambda.\]
	One has a similar expression available with $f_1$ replaced with $\widetilde f_1$. Rearranging $w$ inside the integral completes the proof.
\end{proof}

\subsection{Decomposition for \texorpdfstring{$\op{GL}_2$}{GL(2)}}
\begin{notation}
	Let $\mc H$ be the Hilbert space bundle over $\CC$ whose fiber $\mc H_\lambda$ over $\lambda\in\CC$ consists of functions $f\colon N(\AA_F)P(F)(Z(\AA_F),\omega)\backslash G(\AA_F)\to\CC$ satisfying
	\[f(\op{diag}(u,v)g)=\left|u/v\right|^{\lambda+1/2}\varphi(g),\]
	and satisfying a suitable finiteness condition.
\end{notation}
\begin{remark}
	One can further diagonalize $\mc H_\lambda$ into irreducible representations, but we will frequently not bother to do so.
\end{remark}
\begin{notation}
	Let $\chi\colon[T]\to\CC^\times$ be a character. Define $R_\chi w\colon\op{Ind}_P^G\chi\to\op{Ind}_P^G{^w\chi}$ to be the (restricted) tensor product of the local operators.
\end{notation}
\begin{remark}
	This restricted tensor product is well-defined because it sends the distinguished unramified vector to the distinguished unramified vector when everything is unramified. This is a bounded operator because $R_{\chi_v}w$ has norm one for cofinitely many $v$.
\end{remark}
\begin{remark}
	By expanding out the local calculations, we find $Rw\circ Rw=1$, and $R_\chi w$ is unitary whenever $\chi$ is unitary.
\end{remark}
\begin{theorem}
	The global operator
	\[Mw\colon\int_{X^*(T)_\CC}^\oplus\op{Ind}_P^G\chi\,d\chi\to\int_{X^*(T)_\CC}^\oplus\op{Ind}_P^G{^w\chi}\,d\chi\]
	admits meromorphic continuation, with at worst simple poles located along $s=1/2$, where $\chi_1\chi_2^{-1}=\left|\cdot\right|^s$.
\end{theorem}
\begin{proof}[Sketch]
	Simply define
	\[M_\chi w\coloneqq\frac{L\left(0,\mu\nu^{-1}\right)}{L\left(1,\mu\nu^{-1}\right)\varepsilon\left(0,\mu\nu^{-1}\right)}\cdot R_\chi w,\]
	which is equal to the restricted tensor product of the $M_{\chi_v}w$s.
\end{proof}
\begin{remark}
	Using the global functional equation, one sees that $Mw\circ Mw=1$. Thus, $M_\chi w$ is unitary when $\chi$ is unitary.
\end{remark}
\begin{remark}
	The residue of the pole of $M_\chi w$ at $s=1/2$ can be understood in terms of the inner products
	\[\op{Res}_\chi\left\langle Mw(\widehat f_1),\widehat f_2\right\rangle\simeq\left\langle R_\chi w(\widehat f_1)(\chi),\widehat f_2(\chi)\right\rangle=\left\langle\widehat f_1,\chi\circ\det\right\rangle\left\langle\widehat f_2,\chi\circ\det\right\rangle.\]
	The last equality follows from the local calculation. The relevant constant can be made explicit.
\end{remark}
\begin{theorem}
	One has that $L^2(\op{GL}_2(F)\backslash\op{GL}_2(\AA_F),\omega)$ is isomorphic to
	\[L^2_{\mathrm{cusp}}(\op{GL}_2(F)\backslash\op{GL}_2(\AA_F),\omega)\oplus\widehat{\bigoplus_{\chi_1\chi_2=\omega}}\int_{i\RR/\langle\pm\rangle}^\oplus\op{Ind}_B^G\left(\chi_1\left|\cdot\right|^\lambda,\chi_2\left|\cdot\right|^{-\lambda}\right)\,d\lambda\oplus\widehat{\bigoplus_{\chi^2=\omega}}({\chi\circ\det}).\]
\end{theorem}
\begin{proof}[Sketch]
	The second term is the diagonalization of the (global sections of the) bundle $\mc H$. Anyway, by moving the vertical line of integration, we find that
	\begin{align*}
		\langle F_{f_1},F_{f_2}\rangle &= \frac1{2\pi i}\int_{\op{Re}s=0}\left\langle\widehat f_1(s)+M_sw\widehat f_1(-s),\widehat f_2(-\overline s)\right\rangle\,ds \\
		&\qquad+c\sum_{\chi^2=\omega}\left\langle\widehat f_1(1/2),\chi\circ\det\right\rangle\overline{\left\langle\widehat f_2(1/2),\chi\circ\det\right\rangle},
	\end{align*}
	where the second sum accumulates the relevant poles. Note that $\langle\widehat f_1(1/2),{\chi\circ\det}\rangle$ unfolds into $\langle F_1,{\chi\circ\det}\rangle$.

	We are now in a position to define our isometry.
	\begin{itemize}
		\item Define $a$ to send $C_c^\infty(N(\AA_F)P(F)Z(\AA_F)\backslash G(\AA_F),\omega)$ to the direct integral by sending $f$ to the section $a_f(iy)\coloneqq\frac12(\widehat f(iy)+M_{iy}w(\widehat f(-iy)))$, which we note satisfies $M_{iy}w(a_f(iy))=a_f(-iy)$.
		\item If $\chi^2=\omega$, there is a map $C_c^\infty(N(\AA_F)P(F)Z(\AA_F)\backslash G(\AA_F),\omega)\to\chi\circ\omega$ given by $f\mapsto\langle\widehat f(1/2),{\chi\circ\det}\rangle$.
	\end{itemize}
	The calculation above shows that $F_\bullet\colon C_c^\infty(N(\AA_F)P(F)Z(\AA_F)\backslash G(\AA_F),\omega)\to L^2_{\mathrm{cusp}}(\op{GL}_2(F)\backslash\op{GL}_2(\AA_F),\omega)^\perp$ and the map
	\[\left(a,\widehat\cdot(1/2)\right)\colon C_c^\infty(N(\AA_F)P(F)Z(\AA_F)\backslash G(\AA_F),\omega)\to \widehat{\bigoplus_{\chi_1\chi_2=\omega}}\int_{i\RR/\langle\pm\rangle}^\oplus\op{Ind}_B^G\left(\chi_1\left|\cdot\right|^\lambda,\chi_2\left|\cdot\right|^{-\lambda}\right)\,d\lambda\oplus\widehat{\bigoplus_{\chi^2=\omega}}({\chi\circ\det})\]
	produces the same norm, so their completions are isomorphic. The theorem follows.
\end{proof}

\section{The Spectral Expansion}

\subsection{The Cuspidal Spectrum}
\begin{lemma}
	Choose $f\in C_c^\infty(A_G\backslash G(\AA_F))$, and define $f^*$ by $f^*(g)\coloneqq\overline{f\left(g^{-1}\right)}$. Then $(Rf)^*=R(f^*)$ as operators on $L^2([G])$.
\end{lemma}
\begin{proof}[Sketch]
	Expand
	\[\langle Rf(\varphi),\psi\rangle=\int_{[G]}\int_{A_G\backslash G(\AA_F)}f(g)\varphi(xg)\overline{\psi(x)}\,dg\,dx.\]
	Sending $x\mapsto xg^{-1}$ and then $g\mapsto g^{-1}$ gives the result.
\end{proof}
\begin{proposition}
	Choose a subrepresentation $V\subseteq L^2_{\mathrm{disc}}([G])$. For any $f\in C_c^\infty(A_G\backslash G(\AA_F))$, the operator $Rf|_V$ admits a smooth kernel.
\end{proposition}
\begin{proof}[Sketch]
	Fix a basis $\{\varphi_i\}_i$ of $V$. Then define
	\[K_{f,V}(x,y)\coloneqq\sum_iRf(\varphi_i)(x)\overline{\varphi_i(y)}.\]
	This is the desired kernel once we view $Rf|_V$ as the composite $L^2([G])\onto V\stackrel f\to V\subseteq L^2([G])$. Smoothness is checked by the Dixmier--Malliavin theorem, which lets us reduce to the case where $f=f_1*f_2*f_3$, so $K_{f,V}$ can be viewed as a translate of the group action by $f_1$ and $f_2$.
\end{proof}
\begin{definition}[relative trace]
	Let $H\subseteq G\times G$ be a subgroup for which $H^\circ=H_u\times H_r$; such an $H$ is called \textit{almost reductive}. Define $A_{G,H}\coloneqq H(\AA_F)\cap(A_G\times A_G)$. For a character $\chi\colon A_{G,H}H(F)\backslash H(\AA_F)\to\CC^\times$ and cuspidal $\pi\subseteq L^2_{\mathrm{cusp}}(A_G\backslash G(\AA_F))$ and $f\in C_c^\infty(A_G\backslash G(\AA_F))$, the \textit{relative trace} is
	\[\op{rtr}_\chi\pi(f)\coloneqq\int_{A_{G,H}H(F)\backslash H(\AA_F)}K_{\pi(f)}(h)\chi^{-1}(h)\,dh.\]
\end{definition}
\begin{remark}
	This integral absolutely converges, but it is a little tricky to show. In short, we may pass to the connected $H$ case by a finite-index argument, and then we may pass to the reductive $H$ case because $[H_u]$ is compact. Via Siegel sets, one can basically pass to the case where $H$ is a split torus, and then the result follows upon expanding $K_{\pi(f)}(x,y)=\sum_iRf(\varphi_i)(x)\overline{\varphi_i(y)}$ because functions in $V$ are rapidly decreasing.
\end{remark}
\begin{remark}
	Intuitively, $\op{rtr}_\chi$ is a period against $\chi$. Accordingly, one can show that $\pi\otimes\pi^\lor$ is $\chi$-distin\-guished if and only if the functional $\op{rtr}_\chi\pi\colon C_c^\infty(A_G\backslash G(\AA_F))\to\CC$. The main difficulty lies in upgrading a generic function to a smooth compactly supported one.
\end{remark}
\begin{theorem}[Cuspidal expansion]
	Fix $f\in C_c^\infty(A_G\backslash G(\AA_F))$. Let $H\subseteq G\times G$ be almost reductive, and choose a character $\chi\colon A_{G,H}H(F)\backslash H(\AA_F)\to\CC^\times$. Then for $f\in C_c^\infty(A_G\backslash G(\AA_F))$,
	\[\int_{A_{G,H}H(F)\backslash H(\AA_F)}K_{f,\mathrm{cusp}}(h)\chi^{-1}(h)\,dh=\sum_{\text{cuspidal }\pi}\op{rtr}_\chi\pi(f).\]
\end{theorem}
\begin{proof}[Sketch]
	Once granted absolute convergence, this follows because a basis expansion implies
	\[K_{f,\mathrm{cusp}}=\sum_\pi K_{\pi(f)}\]
	almost everywhere. As usual, by the Dixmier--Malliavin theorem, we may reduce to the case where $f=f_1*f_2*f_3$. Then one finds that
	\[K_{\pi(f)}=(\pi(f_2)\times\pi(\overline {f_1^*}))K_{\pi(f_3)}.\]
	We thus find that $K_{\pi(f)}$ is rapidly decreasing, which is enough to give the required absolute convergence.
\end{proof}
\begin{remark}
	Using the full decomposition of $L^2([G])$, one has a spectral expansion
	\[K_f(x,y)=\sum_{P=MN}n_P^{-1}\sum_{\sigma\subseteq L^2_{\mathrm{disc}}([M])}\int_{i\mf a^{G*}_M}\Bigg(\sum_{\varphi\in\mc B(\sigma)}E(x,I(\sigma,\lambda)(f)(\varphi),\lambda)\overline{E(y,\varphi,\lambda)}\Bigg),\]
	where the sum is over all standard parabolic subgroups, and $\mc B(\sigma)$ is a chosen basis of $\sigma$. Indeed, this is essentially the content of the decomposition of $L^2([G])$ (and its relationship with Eisenstein series).
\end{remark}

\subsection{The Residual Spectrum for \texorpdfstring{$\op{GL}_2$}{GL(2)}}
\begin{definition}[truncation]
	Choose $T>0$. For $\varphi\colon[G]\to\CC$, we define the \textit{truncation}
	\[\land^T\varphi(g)\coloneqq\varphi(g)-\sum_{\gamma\in P(F)\backslash G(F)}\varphi_N(\gamma g)1_{[T,\infty)}(\log H(\gamma g)),\]
	where we recall $H(n\op{diag}(a,b)k)=\left|a/b\right|$. An argument on Siegel sets shows that the sum is locally finite.
\end{definition}
\begin{example}
	If $\varphi$ is cuspidal, the constant term vanishes, so $\land^T\varphi=\varphi$.
\end{example}
\begin{remark}
	For $g\in G(\AA_F)$ satisfying $H(g)>1$, one can show taht any $\gamma\in G(F)$ with $H(\gamma g)>1$ must have $\gamma\in P(F)$. One can use this to show that
	\[\land^T\varphi(g)=\varphi(g)-\varphi_N(g)1_{[T,\infty)}(\log H(g))\]
	if $H(g)>1$.
\end{remark}
\begin{remark}
	Intuitively, the truncation kills constant terms, so it has rapid decrease (e.g., if $\varphi$ is smooth).
\end{remark}
\begin{remark}
	For $T>0$, the operator $\land^T$ is self-adjoint: indeed, a direct expansion and unfolding gives
	\[\left\langle\land^T\varphi_1,\varphi_2\right\rangle=\langle\varphi_1,\varphi_2\rangle-\int_{P(F)N(\AA_F)Z(\AA_F)\backslash G(\AA_F)}\varphi_{1N}(g)\overline{\varphi_{2N}(g)}1_{[T,\infty)}(\log H(g))\,dg.\]
\end{remark}
\begin{remark}
	One can show that $\land^T\land^T\varphi=\land^T\varphi$ by direct calculation. The point is that $\left(\land^T\varphi\right)_N(g)=0$ if $\log H(T)>T$.
\end{remark}
\begin{notation}
	Let $K\in L^2([G]\times[G])$ be a function. Then we define $\land^T_2K(x,y)$ to denote truncation in the second coordinate.
\end{notation}
\begin{example}
	For $f\in C_c^\infty(A_G\backslash G(\AA_F))$, we see that $\land^T_2K_{f,\mathrm{cusp}}(x,y)=K_{f,\mathrm{cusp}}(x,y)$ by direct expansion of the kernel on a basis.
\end{example}
\begin{lemma}
	For $f\in C_c^\infty(A_G\backslash G(\AA_F))$, we have
	\[\lim_{T\to\infty}\int_{[\overline G]}\land^T_2K_{f,\mathrm{res}}(x,x)\,dx=\tr\left(Rf;L^2_{\mathrm{res}}([G])\right).\]
\end{lemma}
\begin{proof}[Sketch]
	Note that $Rf$ is compact on the residual spectrum, so the trace equals
	\[\int_{[\overline G]}K_{f,\mathrm{res}}(x,x)\,dx=\tr\left(Rf;L^2_{\mathrm{res}}([G])\right).\]
	Thus, the difficulty lies in removing the truncation. By expanding on a basis, we see that
	\[K_{f,\mathrm{res}}(x,y)\simeq\sum_{\chi^2=\omega}\Bigg(\chi(\det x)\overline{\chi(\det y)}\int_{\overline G(\AA_F)}f(g)\chi(\det g)\,dg\Bigg).\]
	Truncating in $y$ has the effect of multiplying the $\chi$ summand by $1-\sum_{\gamma}1_{[T,\infty)}(\exp H(\gamma y))$. Thus, one needs to bound
	\[\int_{\overline G(\AA_F)}\Bigg(\sum_{\gamma\in P(F)\backslash G(F)}1_{[T,\infty)}(\exp H(\gamma x))\Bigg)\,dx.\]
	As $T\to\infty$, one shows that this vanishes by unfolding the integral and integrating via the Iwasawa decomposition.
\end{proof}

\subsection{The Continuous Spectrum for \texorpdfstring{$\op{GL}_2$}{GL(2)}}
\begin{notation}
	Following earlier discussion, given a holomorphic section $\varphi$ of $\mc H$, we define the \textit{Eisenstein series} as
	\[E(x,\varphi,\lambda)\coloneqq\sum_{\gamma\in P(F)\backslash G(F)}\varphi_\lambda(\gamma x).\]
	This is an example of a $P$-series.
\end{notation}
\begin{remark}
	By Mellin inversion, for $f\in C_c^\infty(N(\AA_F)(Z(\AA_F),\omega)T(F)\backslash G(\AA_F))$, we see $F_f(g)$ is
	\[\sum_{\gamma\in P(F)\backslash G(F)}f(\gamma g)=\frac1{2\pi i}\sum_{\gamma\in P(F)\backslash G(F)}\int_{\op{Re}\lambda=x}\widehat f_s(\gamma g)\,d\lambda=\frac1{2\pi i}\int_{\op{Re}\lambda=x}E(g,\widehat f,\lambda)\,d\lambda.\]
	Thus, we will be able to use Eisenstein series to discuss the continuous spectrum.
\end{remark}
\begin{theorem}
	Let $\varphi$ be a flat section of $\mc H$. Then $E(x,\varphi,\lambda)$ admits meromorphic continuation to all $\CC$, and its poles agree with the poles of $E_N(x,\varphi,\lambda)$.
\end{theorem}
\begin{proof}[Sketch]
	Note that $E_N(g,\varphi,\lambda)=\varphi_\lambda(g)+M_\lambda w(\varphi_\lambda)(g)$, so it admits meromorphic continuation. For the continuation, we use Whittaker expansion. To this end, we fix a nontrivial character $\psi\colon[N]\to\CC^\times$, and we would like to bound the Fourier coefficients
	\[\int_{[N]}E(ng,\varphi,\lambda)\psi_\alpha(n)^{-1}\,dn\]
	for $\alpha\in F$. If $\alpha=0$, we get the constant term. For $\alpha\in F^\times$, unfolding with the Bruhat decomposition yields that the Fourier coefficient is essentially
	\[\int_{N(\AA_F)}\varphi_\lambda\left(wn\op{diag}(\alpha,1)g\right)\psi(-n)\,dn.\]
	Now, by density, we may as well assume that $\varphi_\lambda$ is a pure tensor and Schwartz--Bruhat, so the integral factors. At finite places $v$, it vanishes if $\left|\alpha\right|_v$ is large enough to keep $\varphi_\lambda$ oscillating enough, so the Fourier coefficients are supported in a fractional ideal of $F$. At the infinite places, we are basically looking at a genuine Fourier coefficient, which is still Schwartz. Convergence follows.
\end{proof}
\begin{corollary}
	Let $\varphi$ be a flat section of $\mc H$. Then $E(x,\varphi,\lambda)$ satisfies the functional equation
	\[E(x,Mw(\varphi),g)=E(x,\varphi,g).\]
\end{corollary}
\begin{proof}[Sketch]
	Both sides have the same constant term, so the difference is cuspidal. But the difference is also a $P$-series, so it is orthogonal to $L^2_{\mathrm{cusp}}([G])$.
\end{proof}
\begin{notation}
	Temporarily, define $S$ to be the projection from $L^2([G])$ to the global sections of $\mc H$, so $S$ factors through the projection to the continuous spectrum.
\end{notation}
\begin{lemma}
	Fix $f\in C_c^\infty((Z(\AA_F),\omega)T(F)N(\AA_F)\backslash G(\AA_F))$ and let $\varphi$ be a flat section of $\mc H$. Then
	\[\langle \varphi,SF_f(iy)\rangle=\frac12\int_{G(F)Z(\AA_F)\backslash G(\AA_F)}E(h_{iy},g)\overline{F_f(g)}\,dg.\]
\end{lemma}
\begin{proof}[Sketch]
	Choose $f_1$ and $f_2$, and define $F_i\coloneqq F_{f_i}$ for brevity. The idea is to compute $\langle F_1,F_2\rangle$ in two ways. On one hand, earlier we computed
	\[\langle F_1,F_2\rangle=\frac1\pi\int_\RR\langle a_{f_1}(iy),a_{f_2}(iy)\rangle\,dy+c\sum_{\chi^2=\omega}\langle F_1,{\chi\circ\det}\rangle\overline{\langle F_2,{\chi\circ\det}\rangle}.\]
	On the other hand, we can expand $F_1$ into an integral of Eisenstein series, which gives
	\[\langle F_1,F_2\rangle=\frac1{2\pi i}\int_{\op{Re}s=x}\left(\int_{\overline G(F)\backslash\overline G(\AA_F)}E(g,\widehat f_1,s)\overline{F_2(g)}\,dg\right)\,ds\]
	for $x>1/2$. Moving the line of integration to $x=0$, we pick up residues from poles, which are the poles of the constant term $E_N$, which all come from the intertwining operator. Thus, we get
	\[\langle F_1,F_2\rangle=\frac1{2\pi}\int_\RR\left(\int_{\overline G(F)\backslash\overline G(\AA_F)}E(g,a_{f_1},iy)\overline{F_2(g)}\,dg\right)\,dy+c\sum_{\chi^2=\omega}\langle F_1,{\chi\circ\det}\rangle\overline{\langle F_2,{\chi\circ\det}\rangle},\]
	where we have implicitly used the functional equation to replace $\widehat f_1$ with $a_{f_1}$. Comparing the two expressions and letting $f_1$ vary completes the proof.
\end{proof}
\begin{proposition}
	Fix $f\in C_c^\infty(A_G\backslash G(\AA_F))$. Let $\{\varphi_\alpha\}$ be a basis of $\mc H_0$, which we extend to flat sections. Then
	\[K_{f,\mathrm{cont}}(g,h)=\frac1{4\pi}\sum_{\alpha,\beta}\int_\RR\langle\pi_{iy}(f)\varphi_\beta,\varphi_\alpha\rangle E(h,\varphi_\alpha,iy)\overline{E(g,\varphi_\beta,iy)}\,dy.\]
\end{proposition}
\begin{proof}[Sketch]
	Note that $S^*S$ just projects $L^2([G])$ onto the continuous spectrum. Thus, after some rearrangement, it is enough to show that
	\[\langle S^*S(Rf)S^*SF_1,F_2\rangle\stackrel?=\iint_{[G]^2}F_1(g)\overline{F_2(h)}\Bigg(\frac1{4\pi}\sum_\alpha\int_\RR E(h,\varphi_\alpha,iy)\overline{E(g,f*\varphi_\alpha,iy)}\,dy\Bigg)\,dg\,dh.\]
	Well, rewrite the left-hand side as $\langle f*SF_1,SF_2\rangle$, and then expand out the inner product in terms of a basis of $\mc H_0$ and apply the lemma.
\end{proof}
\begin{lemma}
	Fix flat sections $\varphi_1$ and $\varphi_2$ of $\mc H$. Then
	\[\left\langle\land^TE(\varphi_1,s),\land^TE(\varphi_2,-\overline s)\right\rangle=2T\langle\varphi_1,\varphi_2\rangle+\langle M_{-s}M'_s\varphi_1,\varphi_2\rangle+\frac1{2s}\left(\langle\varphi_1,M_{-\overline s}\varphi_2\rangle e^{2sT}-\langle M_s\varphi_1,\varphi_2\rangle e^{-2sT}\right).\]
\end{lemma}
\begin{proof}[Sketch]
	Work with generic $s_1,s_2\in\CC$ for a moment. Because we are taking an inner product with a $P$-series, we see
	\[\left\langle E(\varphi_1,s_1),\land^TE(\varphi_2,\overline s_2)\right\rangle\]
	can be expanded into an integral against the constant term of $E(\varphi_1,s_1)$. This constant term is basically $\varphi_1+Mw(\varphi_1)$, so we have brought everything down to integrals of $\varphi_1$ and $\varphi_2$. After using the Iwasawa decomposition to simplify these integrals, we are left with
	\begin{align*}
		&\left(\langle\varphi_1,\varphi_2\rangle e^{(s_1+s_2)T}-\langle M_{s_1}w(\varphi_1),M_{\overline s_2}w(\varphi_2)\rangle e^{-(s_1+s_2)T}\right)\frac1{s_1+s_2} \\
		&\qquad+\left(\langle\varphi_1,M_{\overline s_2}w(\varphi_2)\rangle e^{(s_1-s_2)T}-\langle M_{s_1}w(\varphi_1),\varphi_2\rangle e^{-(s_1-s_2)T}\right)\frac1{s_1-s_2}.
	\end{align*}
	Sending $s_2\to-s_1$ and applying a version of L'H\^opital's rule completes the proof.
\end{proof}
\begin{proposition}
	Fix $f\in C_c^\infty(A_G\backslash G(\AA_F))$. Then
	\begin{align*}
		\int_{[\overline G]}\land_2^TK_{f,\mathrm{cont}}(g,g)\,dg &= T\op{vol}\left(F^\times\backslash\AA_F^1\right)\sum_{\gamma\in\overline M(F)}\int_{M(\AA_F)\backslash G(\AA_F)}f\left(g^{-1}\gamma g\right)\,dg\\
		&\qquad-\frac1{4\pi}\int_\RR\tr M_{-iy}M'_{iy}\pi_{iy}(f)\,dy+\frac14\tr M_0\pi_0(f)+o_T(1).
	\end{align*}
\end{proposition}
\begin{proof}[Sketch]
	This is a very long calculation. Plugging into the expansion of $K_{f,\mathrm{cont}}$, we see that
	\begin{align*}
		\int_{[\overline G]}\land_2^TK_{f,\mathrm{cont}}(g,g)\,dg &= \frac T{2\pi}\sum_{\alpha,\beta}\int_\RR\langle\pi_{iy}(f)\varphi_\beta,\varphi_\alpha\rangle\langle\varphi_\alpha,\varphi_\beta\rangle\,dy \\
		&\qquad-\frac1{4\pi}\sum_{\alpha,\beta}\int_\RR\langle\pi_{iy}(f)\varphi_\beta,\varphi_\alpha\rangle\langle M_{-iy}M'_{iy}\varphi_\alpha,\varphi_\beta\rangle\,dy \\
		&\qquad+\frac1{4\pi}\sum_{\alpha,\beta}\int_\RR\langle\pi_{iy}(f)\varphi_\beta,\varphi_\alpha\rangle\left(\langle\varphi_\alpha,M_{iy}\varphi_\beta\rangle\frac{e^{2Tiy}}{2iy}-\langle M_{iy}\varphi_\alpha,\varphi_\beta\rangle\frac{e^{-2Tiy}}{2iy}\right)\,dy.
	\end{align*}
	The terms are ordered to align with the proposition.
	\begin{itemize}
		\item The first sum rearranges to
		\[\frac T{2\pi}\int_\RR\tr\pi_{iy}(f)\,dy.\]
		One can directly check that the action of $\pi_{iy}(f)$ on $\mc H_{iy}$ is given by the kernel
		\[\int_{N(\AA_F)}\int_{\RR^+}\sum_{\alpha\in F^\times}f\left(x^{-1}n\op{diag}(\alpha at,1)k)\right)\left|t\right|^{-\frac12+iy}\,dn\,d^\times t\]
		as a function of $x$ and the pair $(a,k)\in F^\times\backslash\AA_F^1\times K$. (We need to use the Iwasawa decomposition.) Integrating the kernel along the diagonal gives the first term of the term of the proposition.
		\item Rearrangemnt shows that the second sum corresponds to the second term in the proposition.
		\item The third term rearranges to
		\[\frac1{4\pi}\int_\RR\tr M_{-iy}\pi_{iy}(f)\cdot\frac{e^{2Tiy}-e^{-2Tiy}}{2iy}\,dy+\frac1{4\pi}\int_\RR\frac{\tr M_{-iy}\pi_{iy}(f)-\tr M_{iy}\pi_{iy}(f)}{2iy}\cdot e^{-2Tiy}\,dy.\]
		The right term is a Fourier coefficient, so it goes to zero as $T\to\infty$. The left term is computed using differentiation under the integral sign (by differentiating with respect to $T$).
		\qedhere
	\end{itemize}
\end{proof}

\end{document}